
\newcommand{\compileall}{def}


\documentclass[amstex,twoside]{ctexbook}

% 16开本
\usepackage[
		paperwidth=195mm,
		paperheight=271mm,
		lmargin = 36mm,
		rmargin = 30mm,
		tmargin = 38mm,
		bmargin = 35mm
	]{geometry}

\usepackage{needspace}
\usepackage[colorlinks,linkcolor=black]{hyperref}
\usepackage[below,section]{placeins}
\usepackage{tabularx}
\usepackage{longtable}
\usepackage{multirow}
\usepackage{soul}
\usepackage{amsmath}
\usepackage{wasysym}
\usepackage{amsthm}
\usepackage{stmaryrd}
%\usepackage{eso-pic}


\AddToShipoutPicture{%
\parbox[b][\paperheight]{\paperwidth}{%
	\vfill
\begin{center}
    \includegraphics[width=\textwidth]{pics/draft}\hfill
\end{center}
    \vfill
}}

%=====================================================================================
% font setting
%=====================================================================================
%=====================================================================================
% font setting
%=====================================================================================

\setCJKmonofont[Path="fonts/",BoldFont=AdobeKaitiStd-Regular.otf ]{AdobeFangsongStd-Regular.otf}

\setCJKfamilyfont{zhfs}[Path="fonts/", BoldFont=AdobeKaitiStd-Regular.otf]{AdobeFangsongStd-Regular.otf}


\setCJKmainfont[
		Path="fonts/",
		BoldFont = AdobeHeitiStd-Regular.otf,
		ItalicFont= AdobeKaitiStd-Regular.otf
	]{AdobeSongStd-Light.otf}

\setCJKfamilyfont{KaiTi}[Path="fonts/"]{AdobeKaitiStd-Regular.otf}
\newfontfamily\xingkai[Path="fonts/"]{xinkai.ttf}
\setCJKfamilyfont{xingkai}[Path="fonts/"]{xinkai.ttf}

%=====================================================================================
% define new command and environment 
%=====================================================================================

%=====================================================================================
% define new command and environment 
%=====================================================================================
\newenvironment{notice}{\tt}{}
\newenvironment{insertnote}{ \ttfamily\CJKfamily{KaiTi} }{\vskip 1cm }
\newenvironment{code}{\small\tt\begin{longtable}{p{0.8\textwidth}}}{\end{longtable}}

\newcommand{\RTLpar}{% right-to-left paragraph alignment
  \leftskip=0pt plus .5fil%
  \rightskip=0pt plus -.5fil%
  \parfillskip=0pt plus .5fil%
}

\newenvironment{quotes}[2][0.55]{\pushQED{#2}%
\begin{flushright}%
\begin{minipage}{#1\textwidth}\begin{flushright}\noindent\it\RTLpar}{%
 \\------\popQED{}\end{flushright}\end{minipage}\end{flushright}\vskip 8mm }%

\newcounter{lizi}[chapter]

\newenvironment{example}[1]{ \addtocounter{lizi}{1} \vskip 3ex \bf 例\hskip -0.05em\arabic{chapter}.\arabic{lizi}： #1%
\par %
 \tt\small\CJKfamily{zhfs}}{ \vskip 3ex  }

\newcommand{\theexample}{\arabic{chapter}.\arabic{lizi}}

\makeatletter\newcommand{\chatu}{\@ifstar%
                     \chatuStar%
                     \chatuNoStar%
}\makeatother

\newcommand{\chatuNoStar}[3][scale=0.35]{%
\begin{figure}[h]%
\noindent\centering%
\includegraphics[#1]{pics/#2}%
\caption{#3\label{fig:#2}}%
\end{figure}%
}


\newcommand{\chatuStar}[3][scale=0.35]{%
\begin{figure*}[!h]%
\noindent\centering%
\includegraphics[#1]{pics/#2}%
\label{fig:#2}%
\end{figure*}%
}


\newcommand{\faqref}[1]{~附录A.\nameref{FAQ}~“\nameref{#1}”}
\newcounter{faqs}

\makeatletter
\newcommand\@publisher{}
\newcommand{\publisher}[1]{
\renewcommand\@publisher{#1}
}

\newcommand\@company{}
\newcommand{\company}[1]{
\renewcommand\@company{#1}
}
\makeatother


%=====================================================================================
% style settings
%=====================================================================================

% 阿拉伯数字章节
\setcounter{chapter}{-1}
\CTEXsetup[number={\arabic{chapter}}]{chapter}
% 对齐设置
\tolerance=4500
\punctstyle{hangmobanjiao}
\CJKecglue{}
\CJKsetecglue{} %{\hskip 0.1em plus 0.05em minus 0.05em}

%=====================================================================================
% book info
%=====================================================================================
% book info
%=====================================================================================

\title{Linux枕边书}
\author{微菜}
\company{北京源智天下科技有限公司}
%\publisher{清华大学出版社}
\publisher{XX出版社}

%=====================================================================================


\begin{document}

%=========================================================
% 首页
%=========================================================
{

\thispagestyle{empty}
\makeatletter

\begin{center}


{\fontsize{2cm}{2.5cm}\selectfont \xingkai\CJKfamily{xingkai} \@title

 \vskip 2cm
 
\normalfont  \large  ------ Linux文化入门 
 
 }


\vskip 5cm

 {\LARGE
  \@author ~ 著 }


\end{center}

\null
\vfill



\begin{flushright}
本书由 \LaTeX + \tt{xeCJK} 排版。 \hfill \@publisher ~ 出版
\end{flushright}

\clearpage
\thispagestyle{empty}


\noindent\@author\\
\@title


\noindent Copyright~\textcopyright~2012-\the\year~All rights reversed.

\noindent 本书由 \@company 授权  \@publisher 出版发行。版权所有未经许可，不得已任何凡是复杂或者抄袭本书的任何部分。

\vfill

{
\noindent \bf版权所有，翻印必究
}

\vfill

\noindent 图书在版编目（CIP）数据

XXX

\vfill


\noindent
\hskip -1em
\begin{tabular*}{\textwidth}{>{\bf}ll>{\bf}ll}

出\hfill版\hfill者： &  \@publisher  & 地\hfill址: & 北京清华大学学究大厦  \\
&  \url{http://www.tup.com.cn} &  邮\hfill编: & 100084 \\
&  社总机：010-62770175 &  客户服务： & 010-62776969 \\

责任编辑： & \multicolumn{3}{l}{陈冠军} \\
印\hfill刷\hfill者： & \multicolumn{3}{l}{  XXX 印刷厂  } \\
装\hfill订\hfill者： & \multicolumn{3}{l}{XXX装订厂} \\
发\hfill行\hfill者： & \multicolumn{3}{l}{XXX} \\
开\hfill本： & \multicolumn{3}{l}{ 16开    } \\
版\hfill次: &  \multicolumn{3}{l}{  ~    } \\
书\hfill号: &  \multicolumn{3}{l}{ ISBN XXXX~    } \\
印\hfill数: &  \multicolumn{3}{l}{  ~    } \\
定\hfill价: &  \multicolumn{3}{l}{  ~~~~元    } \\
\hline
\end{tabular*}


\makeatother
\chapter*{致谢!}
\addcontentsline{toc}{chapter}{致谢}

\newpage\thispagestyle{empty}
}

%=========================================================
%目录
%=========================================================
\tableofcontents

%=========================================================
% 序言
%=========================================================
\ifdefined\compileall
	

不知道从什么时候开始用上了Linux，已经不记得刚开始用Linux的时候那段辛苦并乐在其中的学习过程了。
有一天，我重新用起了曾经那么熟悉，那么“简单”的操作系统，结果大失所望。我看到的是一个完全反人类的系统。
它的操作似乎很简单，但是一直总有说不出的感觉，让我觉得它的操作非常别扭。我感觉我第一次接触到Linux的时候也是这种感觉。

有人说，这是一种习惯。而我却更乐意认为那是一种文化。Linux和Windows是一种截然不同的文化。
Linux技术在飞快的进化，各种新技术层出不穷。但是骨子深处的UNIX文化确没变。

有人说，Linux不是UNIX，BSD、Sun Solaris、IBM AIX这些才是UNIX。
没错，这些系统确实继承了AT\&T\footnote{第一个UNIX系统由贝尔实验室开发。贝尔实验室当时属于AT\&T公司。}
的代码\footnote{BSD和AT\&T打了几年官司，后经重写将AT\&T代码完全剔除了。}，还获得了
OpenGroup\footnote{UNIX走向分裂后，几个大公司成立了OpenGroup共同持有UNIX商标。}认证
。但是对我来说，只有Linux才真正在文化上继承了UNIX。
既然Linux的文化源自UNIX，要学习Linux将不得不首先了解UNIX，
了解UNIX从诞生到发展，发展到分裂，分裂到混乱，混乱中统一，统一却失去了UNIX的文化，再到借Linux重生的故事。

故事虽然不是技术，但是了解一个技术诞生的背景，能深刻的理解它。

融入UNIX文化后，再来学习Linux就得心应手了。

但是我不愿意独立一个章节介绍UNIX历史，因为这本书是带领大家学习Linux的，不是故事书，
所以我把整书的序部分用来介绍UNIX的历史文化。因为序比较长，干脆就叫第\thechapter章。

\section{为什么又写一本书}

我不记得自己为了学习Linux查阅了多少资料，翻烂了多少本书。最终翻了那么多书后发现，
现在讲解Linux入门书籍多数都讲求“速成”而采取填鸭式讲解。就算速成后成功学会Linux，最终也会因为不了解Linux的文化而处处碰壁。
然后到网上发牢骚抱怨Linux不人性化。

\begin{notice}
习惯于Windows的人无法习惯于Linux的操作。所有与Windows不一致的操作都被认为“不人性化”。
\end{notice}

必须首先要让读者了解到UNIX文化的精髓，而后开始学习Linux，抛开一切对Windows的比较。这样学会的Linux，遇事首先会以UNIX的思维模式思考，然后提出非常UNIXy\footnote{UNIX的，UNIX化的。}的解决办法。

所以按照人对事物的认知，我对Linux的教学过程的安排将不同于其他的书，我的重点在“文化”，介绍原理而不简单的介绍的用法。
确实我会在书中有专门的章节安排详细讲解个各种实用工具，包括服务器的搭建，但读者可以把它当作是一个手册，用到的时候再翻阅。

在对发行版的考量上，我采用了Gentoo作为本书的参考发行版。在市场上还没有哪本入门书是使用Gentoo作为发行版的。
对于使用人数最多的Ubuntu\footnote{非常流行的一个发行版}发行版，我想一定存在非常非常多的教程。但是没有哪个教程能让我满意。
不是介绍的太过于片面而没有考虑到其他发行版，就是过分介绍了发行版之间的不同点。
原因就是Ubuntu（还有Fedora\footnote{RedHat的社区发行版，其改动将回馈到RedHat企业版。}这些）发行版，并不是只将软件进行打包，还进行了大量的修改使之符合发行版的风格。

\begin{notice}
就拿 /etc/init.d/ 下面的服务启动脚本来说吧，没有哪个发行版的脚本是一致的。
都经过了大量的修改。只介绍一个发行版还不够，但是为了“全面”而大而全式的把所有的发行版都介绍一遍，对初学者来说简直就是个噩梦。他会发现他需要为每个发行版都学习一次init脚本。
\end{notice}

如果说哪个发行版最尊重上游软件\footnote{所谓上游，就是软件的作者。发行版将软件从原始站点下载后重新打包，如果称发行版为下游，那原始软件开发者就是上游。}，进行最小的改动的话，那就是Gentoo了。

Gentoo几乎所有的软件包都是直接从上游代码编译而来，中间很少有加自己的补丁\footnote{对一个软件的修改被形象的称呼为打补丁。}。
有限的几个补丁也是为了修复bug而不是对软件进行修改。

这种接近上游软件的状态非常适合初学者的学习。

\section{UNIX的史前文明}

当你看《天宫韵音》\footnote{86版本的西游记电视剧的片头曲}的时候，最先播放的画面就是悟空出世。第一个音符响起的时候，悟空从绷床\footnote{没错！仔细看，确实是绷床！后期处理居然没处理掉。BUG无处不在啊！}上腾空而起。悟空就这样开天辟地般高调出世了。

但是UNIX的出世却非常低调，低调到需要介绍他那高调的前身\textbf{MULTICS}。

Multics的全称是\textit{MULTiplexed Information and Computing System}\footnote{注音这里的大小写混合。这种混合方式旨在告诉读者缩写的由来。}。我们先来了解一下计算机早期的背景。



\begin{insertnote}
\subsection*{Multics背景}

\chatu{TwoWomenOperatingENIAC}{为世界上的第一台电脑编程\protect\footnotemark}%
	\footnotetext{图片来自维基百科~\url{http://en.wikipedia.org/wiki/File:Two_women_operating_ENIAC.gif}}


计算机诞生的时候并没有操作系统，甚至连程序都没有。操作员重新插拔系统电路连接来实现某个功能。
如图\thefigure那样两个人在重新插拔电路为世界上的第一台电脑编程。

而今天的电脑只要执行不同的程序就能有不同的功能。程序的发明是个非常重要的进步\footnote{通过执行程序而不是直接修改电路实现计算功能，程序和数据一同存储，执行的时候同时进入主存储器，这种结构的计算机最早由冯.诺曼提出，又叫冯.诺曼结构。我们现在的个人电脑就是冯.诺曼结构的。}。有了程序就可以预先编撰好程序，只要输入了程序就能执行了，而不需要费时的重新插拔电路。

那个时候计算机程序的输入靠的是穿孔纸\footnote{用纸啊！太不环保了！}。一长条纸穿进去，程序就输入了，然后等待执行完毕。执行完毕后将结果打印到纸上
\footnote{别想多了，是简单的机械打印机。只能打印0\textasciitilde{}9数字、26个字母和特殊符号。有时就是这点功能都不能保证。可能只能打大写的A-Z而已。UNIX少说多做的传统就来源于此，不仅是为了节约纸张，更是节约打印时间。}。
每次程序执行完毕后操作员要准备好下一叠纸，按下机器的复位按钮机器重新开始读取纸张并执行....周而复始。

如果程序的执行不需要人准备下一叠纸就好了。第一个操作系统就在这样的需求下诞生了。不用人的帮助计算机就能批量执行程序。尽管一次还是只执行一个程序。但是能做到前一个程序执行完毕后自动读取下一个程序执行。

有时候一个程序要执行很长的时间，执行完毕后要花更长的时间打印结果，这个时候处理器的资源就是闲置了\footnote{写小看这个闲置，那个时候的计算机比劳斯莱斯贵多了。}。如果前一个程序在打印结果的时候后一个程序能开始执行就好了，这样就能最大限度的利用处理器的资源。

于是最初的多任务操作系统应运而生。它能将处理器的时间分配给多个程序，虽然在任意瞬间还是只有一个程序在执行，但是由于快速的在不同的程序间进行切换\footnote{专业术语叫上下文切换（Context Switch），所谓上下文切换，就是将当前程序的状态保存起来，然后调入先前保持的另一个程序的状态。这样那个被调入的程序就可以从被切换打断的那个地方继续执行。}，看起来就好像多个程序同时执行了。如图 \ref{fig:ContextSwitch}  所示。

\chatu[scale=0.65]{ContextSwitch}{上下文切换}

但是那个时候的计算机，还不能进行实时交互。程序员写好纸条后，要交给计算机管理员\footnote{确切的说是操作员}，管理员将纸条输入计算机，然后等待计算机将执行结果输出。最后输出结果交给程序员，这样程序员才能知道程序的执行结果。


随着计算机性能的提高，出现了一种分时操作系统（Time Sharing Operating System）。 
它可以让大型主机提供数个终端，用户通过终端使用主机。
终端经一条很长很长的线缆\footnote{也可以通过电话线呢！}连接到主机上\footnote{那个时候主机很大，要放到专门的地方嘛。}。
当时的终端就是一个电传打字机和一个机械打印机，只具有输入/输出的功能，本身完全不具任何运算或者软件安装的能力。要使用计算机的人将程序或者命令通过终端输给主机，主机将执行结果传回终端。最后来屏幕代替了打印机，成为视频终端。

\chatu[scale=0.8]{ctss}{分时操作系统}

如此一来，无论主机在哪里，只要在终端機前面进行输入输出的作业，就可利用主机提供的功能了。

前面已经讲过，之前的系统只能由一个操作员去操作。而分时系统为每一个使用计算机的人提供了一个终端。比较先进的主机大概能同时为30个终端服务。 至此，计算机终于能实时交互\footnote{其实只是接近实时。不过已经是一种进步了，不是么？}了。

Multics就是在分时系统的基础上开发的。

\end{insertnote}

MIT\footnote{麻省理工大学}计算机研究中心研究出来的CTSS(Compatible Time Sharing System)是世界上最早的分时系统，它在1961年首次演示。
正如一切软件开发者所遵循的鉄律：来不及添加的功能就放到第二版做吧！Multics就是CTSS来不及添加更多功能后希望开发的全功能第二版。1964年，MIT、通用电器（General Electric）和贝尔实验室\footnote{根据Multics Myths网站（一个网站致力于破除有关Multics的谣言）的说法，贝尔实验室应该是在1965年的时候才参加的。
}通力合作，计划开发一个完美的CTSS第二代。

Multics计划的目的是想要让大型主机可以\textbf{}为超过300个以上的终端服务和一整套完整的信息处理系统。
Multics支持模块化的硬件配置：CPU、内存，磁盘等硬件都可以在需要的时候在系统运行的过程中被动态增加或移除。在当时的系统差不多都是由汇编语言写成的时候，Multics就大胆的提出了使用高级语言\footnote{当时还没C语言呢！}编写整个系统。

Multics先进的系统模型的结果就是：从1964年立项1969年，整整5年过去了，Multics的开发一再受挫。
Multics迟迟不能完成，预定的目标越来越难以达成，在投入了巨大的人力物力资源后，
贝尔实验室终于决定在这场注定失败的项目\footnote{一种说法是Multics最终达成了设定的目标，所以不能算是一个失败的项目。但是我且以为Multics没能在期望的时间内达到设定的目标，就是失败的。}中提前退出。

但是Multics激发的灵感------高级语言编写的内核、文件系统、终端实时交互、多任务、多用户等------却在今后几十年内被延续到各种各样的系统上。这些最终将在接下来的几十年时间里分别在不同的系统上得以实现。

Multics是个优秀的系统。在参加Multics项目的日子里，贝尔实验室的研究人员已经习惯了不用等待就能通过交互式终端使用大型电脑的功能。
但是退出Multics项目后，他们又得继续原始生活了。由于习惯了Multics带来的便利，贝尔实验室的研究人员将发明一个简化的Multics，这个发明就是\textbf{UNIX}。

\section{UNIX的故事}

很显然Ken Thompson不愿意放弃Multics提供的交互式终端，
而且为了能继续玩他在Multics项目做研究的时候编写的“星际旅行”游戏，他头脑风暴了一把，用实验室废弃的一台PDP-7写了一个简陋的操作系统。
Ken Thompson为PDP-7开发的小型系统被称呼为UNICS（UNiplexed Information and Computing System） 。 这个名字的由来是因为Ken Thompson认为Multics太复杂，而Multics的功能可以用更简单的方式去实现。Dennis Ritchie注意到了他的同事为PDP-7开发的简单系统。两个人亲密无间的合作就这样开始了。

\chatu[scale=0.2]{KandRwithPDP7}{Dennis Ritchie（站）和Ken Thompson（坐）在工作}

Dennis Ritchie和Dennis Ritchie的加入为Ken Thompson的系统带来的第一个变化就是高级语言。他们开始计划用B语言\footnote{C语言的前身，Ken Thompson创造的。}重写这个系统。

在重写UNICS的过程中，Dennis Ritchie改进了B语言，使之成为更加实用的C语言。Dennis Ritchie和Ken Thompson一起实现了C语言，并在1973年用C语言重写了UNICS。哦，不对，是UNIX。因为谐音的关系，不知道什么时候开始UNICS就被称呼为UNIX了。

UNIX很快就在贝尔实验室流行开来了。
1974年，Dennis和Ken第一次对外公布了UNIX。他们在《美国计算机通信》(Communications of the ACM)上发表了介绍UNIX的论文。
这篇论文使得外界对UNIX产生了极大的兴趣，纷纷开始向贝尔实验室索要UNIX。

老实说，AT\&T是个大公司。恩不对，还是个垄断公司。AT\&T垄断了美国的电话网。
根据1958年为解决反托拉斯案例达成的和解协议，AT\&T被禁止进入计算机相关的商业领域。
所以,Unix不能够成为一种商品。实际上，根据和解协议的规定，贝尔实验室必须将非电话业务的技术许可给任何提出要求的人。
Ken Thompson开始默默回应那些请求，将包含着UNIX源代码的磁带和磁盘一包包地寄送出去——据传说，每包里都有一张字条,写
着“love，ken”（爱你的，ken）。

UNIX的源代码在高校间广泛传播\footnote{UNIX因此获得了一个非常重要的传统“附带源代码”。正是这个传统催生了后来的开源运动。}，
其中有一份代码被寄到了加州大学伯克利分校。在那里，UNIX被孵化，催生了BSD（Berkeley Software Distribution）。

在伯克利，一些重要的UNIX功能被开发出来。最重要的莫过于socket了。在BSD开发出socket之前，UNIX的网络功能简陋得可怜\footnote{UNIX有个uucp命令可以通过电话线执行远程拷贝，仅此而已。}。


AT\&T对UNIX源代码的“慷慨\footnote{没那么好心。不过是AT\&T还没意识到UNIX的商业价值，同时暂时性的被反托拉斯和解协议束缚了。}”奉献促成了UNIX的繁荣。

大量的新功能被开发出来，各种功能补丁雪花般飘到贝尔实验室。各种版本的UNIX不断涌现，商业化的UNIX系统也出现了。
伯克利开发BSD的Bill Joy加入了新成立的Sun\footnote{是Stanford University Network的缩写。2011年被甲骨文（Oracle）收购。}公司开发基于BSD的SunOS\footnote{SunOS从 5.0版本以后就不再基于BSD开发了。}。
就连微软也曾开发名叫Xenix的UNIX系统\footnote{但是微软把UNIX作为一个产品的热情并没有持续多久。}。


\subsection{UNIX混乱和分裂}

UNIX世界百花争鸣的日子随着AT\&T公司的拆分一去不复返。1983年，美国司法部对在针对AT\&T的第二起反托拉斯诉讼中获胜，AT\&T和贝尔实验室被拆分。这次判决将AT\&T从1958年的禁止将UNIX产品化的和解协议中解脱了出来。AT\&T马上忙不迭地将UNIX商业化------这一举措差点扼杀了UNIX。

那时，没有人意识到，UNIX的产业化会破坏UNIX源码的自由交流，而恰是后者滋养了UNIX系统早期的活力。AT\&T收回了原先的UNIX源码的许可，挂上高价（而在这之前，AT\&T只对UNIX源码收取象征性的费用），对源码散发严加防护。来自高校的贡献开始枯竭。

刚刚进入UNIX市场的几家大公司马上犯下了最重大的战略错误------试图通过差异化寻求有力地位------这个策略导致了
UNIX的分歧，抛弃了UNIX的跨平台兼容性,造成了UNIX市场分割。

AT\&T 拆分后的数年内，UNIX社区忙着UNIX大战的第一阶段-----System V Unix（AT\&T的将UNIX商业化后推出的UNIX）和 BSD UNIX之间的内部争吵。争吵分成不同的层面，有些属于技术层面（socket vs stream,BSD tty 对 System V termio），有些则属于文化层面。
AT\&T对UNIX的商业化使得UNIX作为统一的操作系统环境分崩离析。市场上充斥着各种互不兼容的UNIX。
给软件开发和系统维护带来了沉重的负担。

\subsection{统一的UNIX和POSIX标准}

1983年/usr/group标准组开始了以调解System V UNIX和BSD UNIX的API为目标的严肃的标准化工作。
随后IEEE以/usr/group标准为基础制定了POSIX
\footnote{其正式称呼为IEEE 1003，而国际标准名称为ISO/IEC9945。POSIX这个名称是由理查德·斯托曼应IEEE的要求而提议的一个易于记忆的名称。}
（Portable Operating System Interface\footnote{最初这个X是 of UNIX最后的那个X。但是后来鉴于POSIX可以用于非UNIX（这些系统被称为UNIX-like，如Linux和Minix）系统，于是去掉了of UNIX但是保留了X。 }，可移植操作系统接口）标准。并于1985年发布了第一版。

POSIX标准定义了UNIX（和UNIX-like）系统的API\footnote{ 参考\faqref{API}}（Application Programing Interface）和shell命令所要遵循的标准。

后续开发的各种UNIX版本也严格遵循这个标准。

\begin{notice}
注意：POSIX标准是源码级兼容的标准。意思是在一个系统上编写的软件可以在另一个系统上“无修改的”通过编译并正常运行。
但是如果是编译好的可执行程序拿到另一个UNIX系统下“不保证”能执行。

但是对于shell脚本就另当别论了。POSIX定义了一个shell所必须具备的功能和行为。只使用POSIX定义的shell功能的脚本可以未经修改直接在另一个UNIX系统上执行。
\end{notice}

那么UNIX的故事到这里就告一段落了。Linux文化的的另一个重要来源GNU也要登上历史舞台了。

\section{GNU 的故事}

GNU也不是什么开天辟地的存在。早在GNU出现之前，类似GNU提倡的那种开放源代码的黑客精神就已经存在了。他们相互共享软件的代码，相互修改对方的软件。相信软件是一种自由的可以共享的东西。他们当中最杰出的一名，在黑客文化行将灭亡的时候，发出了著名的《GNU宣言》。这个人就是Richard Matthew Stallman\footnote{他的简称RMS比他的名字更出名}。

\chatu{rms-full-size}{Richard Matthew Stallman}

\subsection{认识自由的含义}

1971年RMS进入哈佛大学学习，同年受聘于麻省理工学院人工智能实验室（AI Laboratory），成为一名职业黑客。

RMS在人工智能实验室的最出名的发明就是Emacs编辑器。Emacs打从一开始就以源码的形式进行发布。Emacs出色的设计使得它非常的流行。许多人对Emacs进行了改进。这些改进最终都回到了RMS的手里。这种共享代码的方式深受黑客们的喜爱。对于RMS来说，自由的共享软件的源代码是多么自然的事情。

但是商业化的浪潮也侵袭到了MIT。RMS感觉到他熟悉的黑客圈开始消失。

\begin{insertnote}
\subsection*{该死的私有代码}
在软件私有化的浪潮开始前，所有的软件销售的时候都是附带源码的（不过修改代码重新发布的权利还是没有的）。软件通常是作为硬件的“附属品”。
人工智能实验室购买了许多的硬件设备。当硬件驱动有错误的时候，RMS和他的黑客同事们总是能修正这些硬件驱动的BUG。他们甚至能开发出新的驱动让硬件提供不同的功能。

后来开发商逐渐的都不再提供驱动软件的源代码了。这样驱动程序出了问题这些黑客们只能对着机器发呆等待厂家修正。
可想而知RMS心里有多么的痛恨这些私有代码。如果他们有这些糟糕的软件的代码，就能修改它们而不是被动的等待厂商的修复。
\end{insertnote}

进入80年代，他的同事们许多都离开了实验室开始创立公司。
RMS觉得自己被逐出了黑客的伊甸园,他把这一切都归咎于私有代码。

\subsection{为了自由}

RMS认为所有的软件都应该是自由的。自由的意思就是:
\begin{itemize}
\item 运行的自由。你可以自由的决定运行或者不运行这个软件。运行的话怎么运行这个软件，在何种环境下运行\footnote{可别小看这个自由，许多软件规定只能在家运行，不能在作为商业用途在公司运行。}。
\item 修改的自由。你可以获得程序的源代码，并进行修改。
\item 再发布的自由。你可以自由的传播不论是你修改过的还是未修改过的软件。
\end{itemize}

RMS可以编写更多的自由的软件来做抗争。但是只要一个系统的核心------操作系统------是私有的软件，写再多的自由应用程序又有何意义呢？
于是RMS要做的第一件事情就是创建一个自由的操作系统。UNIX本来是自由的，却被AT\&T变成了私有的软件。于是RMS想到编写一个自由的UNIX兼容系统。这个项目的名字就是GNU，是Gnu is Not Unix\footnote{这种递归缩写在很多项目里用到。}的递归缩写。

造一个自由的操作系统首先需要编译器。于是他开始开发C语言的编译器。也就是GCC。到了1985年，RMS发现靠一个人的力量做事情太慢了，他需要帮手。于是他成立了自由软件基金（Free Software Foundation 缩写FSF）。同年RMS撰写了著名的自由宣言《GNU宣言》（The GNU Manifesto）。

1986年，GDB\footnote{GNU Debuger，GNU调式器}发布。次年，第一版GCC\footnote{GNU C Compiler，GNU C编译器。后来因为GCC能支持其他语言，GCC被改名为 GNU Compiler Collection，GNU编译器套件} 发布。至此开发一个软件所需要的3大件------编辑器、调式器、编译器------齐备。

RMS面临的另外一个问题是，他需要保护自由软件的成果不至于被偷窃。UNIX就是一个活生生的由共享源码的自由的软件变成了AT\&T的商业软件的例子。于是在1989年，RMS编写了著名的GPL\footnote{GNU Public License，GNU公共许可协议}协议。

\subsection{以Copyright保护Copyleft}
%\subsection{只有自由软件的完整操作系统}

自由软件的目的是\textbf{去版权化},也被 称为Copyleft。
但是首先要用版权保护自己。软件的作者通常会随软件附带一份许可协议，这份许可协议由版权法保护，用户同意即生效。而且要使用该软件用户必须同意（当然前提是这个国家对版权的保护非常好。）。
用户可以同意并使用该软件或者不同意并退回所购买的软件\footnote{这有点像是霸王条款？}。通常这样的许可协议是用来保护开发者的：微软公司就在他的许可条款里堂而皇之的说因使用微软的软件导致的损失微软免责，这种免责条款被大多数公司所使用。

GPL却利用了这点来\textbf{保证用户的权利}。用户被赋予了和开发者同等的权利，只要保证一点：他同意继续保持软件为GPL授权。

\begin{itemize}
\item 当你下载了GCC，你有自由执行它的权利。GPL授权的GCC不会在许可协议里说，“嘿伙计，你只能在Windows下执行该软件\footnote{微软的Office软件有此条款。}”

\item 当你下载了GCC，你就拥有了GCC开发者相同的权利。你可以查看GCC的代码------FSF的网站上提供GCC源码下载------就和GCC开发者能查看GCC的代码一样。

\item  当你下载了GCC，你就拥有了GCC开发者相同的权利。你可以修改GCC的代码------就和GCC开发者能修改GCC的代码一样。

\item  当你下载了GCC，你就拥有了GCC开发者相同的权利。你可以把GCC修改后（或者没有修改）分享给任何一个人------就和GCC开发者能把这个软件发给你一样，而不会在许可协议里说，“嘿伙计，你不能把本软件拷贝给任何一个人，否则就是非法拷贝，是盗版”。

\end{itemize}

所有的这些权利你都能获得，只要你同意
\begin{itemize}
\item 当你再发行这个修改或者未修改后的软件的时候，你保证不修改它的许可协议。让它的GPL协议传播给下一代。
\item 当你在自己的软件里用到了一个GPL授权的软件的一部分或者全部的代码的时候------就算你没有将这个软件的代码直接包含到你的软件中，而只是使用里它提供的函数库------你保证将你的软件同样以GPL协议授权。
\end{itemize}

GPL就是通过如上的传染性保证自己能像病毒一样迅速传播。随着GPL授权的软件愈来愈多，最后几乎不可能编写一个完全不依赖GPL代码的软件了。这样所有的软件最终都将被迫GPL化\footnote{我不知道这个词语之前有没有人用过，我希望读者能理解这里个词语的含义，就是改变原先的许可协议为GPL。}。

到1992年，GNU项目开发出了一系列的软件，包括编辑器、开发工具链（编译器，调试器，C库，autotools\footnote{参考}）、shell和系统工具等等，再有一个内核就是完整的自由的操作系统了。但是GNU的Hurd内核一再跳票，距离发布的日子遥遥无期。这个时候GPL授权的Linux内核刚好填补了Hurd的位置，使GNU成了一个彻底完整的系统。至此UNIX的文化经过80年代的短暂中断后在GNU和Linux的组合系统中延续。

\else
	\addtocounter{chapter}{1}
\fi


%=========================================================
% 正文开始
%=========================================================

\ifdefined\compileall


UNIX的哲学到底是什么？UNIX的哲学一言以蔽之，就是：\chatu*{kiss}{KISS}


UNIX无处不在，生命力顽强。从1969年诞生的那一刻起，没有人想到它会如此成功。得益于C语言，UNIX以其绝佳的可移植性轻松的被移植到各种类型的机器上。小到各种移动设备，达到超级计算机，UNIX都能在这些设备上畅快的运行。恐怕没有哪个操作系统能像UNIX那样拥有持续40多年的生命力。UNIX不仅存在于UNIX操作系统本身。许多非UNIX系统都能找到UNIX的影子。


%TODO:  主要的系统管理思想和一些小工具也在这里介绍，比如 ps grep 用户和组，用户权限的概念也在这里讲解

\section{一切都是文件}

这是最明显的一条UNIX设计哲学了。通过将一切设备抽象为文件，就可以使用统一的方法去操作设备。不管设备是磁带机也好，猫也罢，都使用同样文件操作。唯一不同的，也只有操作的设备文件名不一样而已。彻头彻尾的反UNIX，就是将设备和文件系统隔离。

\begin{example}{备份 MBR}

在 Linux下，第一个磁盘的设备文件是 /dev/sda , 第二个磁盘的设备文件是 /dev/sdb 
要备份第一个硬盘的 MBR ,可以使用  dd 命令。

\begin{code}
dd if=/dev/sda of=mbr.bak bs=512 count=1
\end{code}

使用这条命令就可以把硬盘的第一个扇区备份到 mbr.bak 文件了。对于其他的UNIX系统，只有对硬盘设备的命名不一致而已，没有其他区别。

而彻头彻尾的反UNIX，就需要使用特殊的备份MBR/分区表的“工具”。
\end{example}


\begin{example}{查看硬扇区}

直接查看硬盘的扇区，就和查看一般文件没有任何区别。使用二进制编辑器打开磁盘设备就可以了。
如果只是为了查看一下，还有个非常方便的工具：\textbf{hexdump}

\begin{code}
hexdump /dev/sda
\end{code}

这条命令就可以把第一块硬盘的扇区数据直接显示到屏幕上，按Ctrl-C终止。想屏幕不要快速滚动的话，请仔细阅读下面的章节，你会找到办法的。

而彻头彻尾的反UNIX，就需要使用特殊的查看磁盘扇区的“工具”。
\end{example}

从两个例子可以看出，UNIX将硬盘看作一个普通的文件（只不过文件大小固定为磁盘容量）。像这样的设备文件UNIX将他们组织起来统一放到/dev目录下面。
/dev 目录下的文件一度是由内核进行管理的，不过Linux后来放弃了内核管理/dev的办法而改为由一个叫udev\footnote{udev最后被systemd合并。}的程序进行管理。


\section{鼓励程序间协作}

UNIX鼓励短小精湛的程序，不鼓励编写大而全式的程序\footnote{据说是因为当年PDP-7内存有限}。正是因为UNIX鼓励短小的程序，并相互配合，所以创造了无穷的灵活性。
UNIX通过组合各种小程序，能发挥出各个小工具独立工作的时候不能完成的功能。

实现程序间协作的最伟大的发明就是\textbf{管道}。通过管道，一个程序的输出可以变成另一个程序的输入，将多个程序进行串街------每个程序变成这条处理链上的一个\textbf{过滤器}------组合的这些工具能产生工具的开发者都没有预料到的功能。

\chatu{Douglas-McIlroy}{管道发明人 Malcolm Douglas McIlroy}

图\thefigure 为管道发明人Malcolm Douglas McIlroy。 \chatu[width=\textwidth]{pipe-draft}{管道} 图\thefigure 为他发明管道的时候提出的建议的打字稿扫描图片。

管道能够让你用简单的程序去构 造更复杂的程序。管道能够以意想不到的方式去使用命令,管道使得你可以用更简单的方式去实现复杂的功能。
UNIX简单就是美的设计哲学，让管道这种强大的功能简单到只需要使用“|”就可以实现了。

\begin{example}{还是查看硬扇区}

前面说过，直接查看硬盘的扇区，只要执行如下的指令：
\begin{code}
hexdump /dev/sda
\end{code}

想屏幕不要快速滚动的话，并不需要修改 hexdump 程序添加滚屏功能。这个功能只需要交给一个专门为\textbf{滚屏}而开发的程序: less。

只要执行这样的命令就能实现一次查看一个屏幕的内容。还可以向上回滚。
\begin{code}
hexdump /dev/sda | less
\end{code}

less 的功能就是将自己的输入以一个屏幕的内容显示，超过屏幕能显示范围的内容要等用户使用滚屏命令（PageUp PageDown 空格键和回车键）让less滚动到用户想要观察的位置。

如果我不太喜欢 less 提供的功能，我还有more呢！
\begin{code}
hexdump /dev/sda | more
\end{code}

more和less在功能上大同小异。

而彻头彻尾的反UNIX，就需要给所有有需要的程序添加滚动屏幕的支持。

\end{example}

UNIX鼓励程序之间通过管道或者其他IPC（Inter-Process Communication 进程间通信）进行协作。这样一个程序就可以专注于本质工作，不需要为了“多功能”而开发庞大复杂的软件。这就引出了下一个我要讲的UNIX哲学。

\section{一个程序只做一件事情并把它做好}

UNIX鼓励程序间的协作，一个程序只需要专注于本职的工作，并做到最好。UNIX下有各种简单的工具。但是这些简单的工具却将自己的本质工作做到了最好。
所以UNIX下通常不会有功能繁多，大而全式的软件。但是你有非常丰富的功能单一却很强大的小工具，并拥有将工具进行组合的能力。

而彻头彻尾的反UNIX，让每个软件都倾向于包含你可能用到的一切功能，试图提供一个“家长式”的程序。

\begin{example}{下载工具wget}

wget 是一个命令行下的下载工具。wget功能非常明确，http/ftp 下载。使用起来也非常简单
\begin{code}
wget [url]
\end{code}
但是wget做到了一件事——下载——并将它做最好。
但是，如果给你的是一个下载文件列表，wget可不会去支持它。这是和下载无关的功能。但是，你能通过组合命令实现：
\begin{code}
cat urllist | xargs wget
\end{code}
在这里，我们又利用了一个名字叫xargs的程序，xargs的程序功能很简单，就是接受标准输入，为每个输入的字符串作为最后一个参数调用一次指定的命令。
cat 是个功能非常简单的程序，将文件打印到标准输出。这个标准输出被\underline{\textbf{管道}}重定向给了xargs做为标准输入。

这个命令组合后的结果就是：为urllist里列出的每个url调用一次wget进行下载。自然也就是为wget实现了“下载文件列表里的url”的功能了。
\end{example}

UNIX就是因为提供了管道这种强大到底协作工具，所以得以鼓励开发各种功能单一的精湛的小程序。而功能单一的程序大大简化了软件的开发，也使得UNIX下的软件质量非常高。

\section{使用文本流，因为这个是通用的协议}

如果说，UNIX鼓励编写功能简单的程序，必定是因为它提供了程序间协作的机制，而要利用这样的机制，程序间唯有使用统一的协议才能实现正确的协作。
这个统一的协议就是文本流。

UNIX反对使用专有的二进制格式。UNIX天然的支持流格式\footnote{一个接着一个的字节}，连文件都被看成是无格式的字节流。
因为输出不仅可以被管道送给另一个程序处理，更有可能会直接打印到终端。
所以对于输入输出，UNIX更是鼓励文本字节流，也就是只使用可打印的\footnote{26个大小写字母，0到9阿拉伯数字，标点符号，空格和换行}文本字母。并且将这种文本用于进程间的通信。

\begin{example}{文本协议HTTP}
目前广泛使用的HTTP协议就是一个文本协议的例子。

比如你访问 \url{http://www.example.com/example.html} ， 浏览器连接到 www.example.com 主机的80端口，发出一个
\begin{code}
GET /example.html HTTP/1.1
\end{code}
请求。服务器收到后，如果找到了example.html文件，就返回
\begin{code}
HTTP/1.1 200 OK
\end{code}
并空1行后紧接着输出example.html文件的内容。

要模拟浏览器的请求，可以使用\textbf{telnet}这个程序。
\begin{code}
\#telnet  www.example.com 80\\
Trying \st{***.***.***.***}...\\
Connected to www.example.com.\\
Escape character is '\textasciicircum'.{ <--  \it 以上为telnet输出 }\\
GET /example.html HTTP/1.1 { <--  \it 自己输入的，输入后敲两下回车 }

\\

HTTP/1.1 200 OK\\
Date: Wed, 19 Sep 2012 15:41:40 GMT\\
Content-Length: 31\\
Content-Type: text/html; {  <-- \it 下面有一行空行}

\\

<html>hello example.com</html>

Connection closed by foreign host. <-- \it telnet输出
\end{code}

可以看到，HTTP协议使用文本模式进行通信。服务器的回应也是文本流。而且是行模式，没条信息存放一行。如果有多个信息，可以多行存放。HTTP协议没有规HTTP头的大小，HTTP头和传输的内容之间使用空行隔开\footnote{这就要求HTTP头本身没有空行。}。

\end{example}

除了HTTP，互联网上常用的协议，比如FTP、SMTP\footnote{发送邮件用的}、POP3\footnote{用于客户端从服务器接收邮件}都是基于文本的协议。
HTTP这种只有HTTP头是文本，后续发送的文件内容（尤其是二进制文件来说）是字节流。但是SMTP全程连接使用的都是文本。对于发送不可打印的字符（比如以一张图片作为附件），SMTP需要将二进制字节流转化为文本格式。这是文本协议的一个缺点。

但是瑕不掩玉，文本协议的好处有很多，比如可以无限扩展。

还是拿 HTTP协议来说（方便起见，我没有设计HTTP请求的二进制版本。）。刚刚访问 www.example.com获了example.html。该文件有31个字节长度，HTTP头中的信息Content-Length: 31也告诉了我们文件大小。如果是二进制协议，一般会设计使用一个32bit的整数代表返回的文件大小。如果需要处理很大的文件，协议恐怕需要使用64bit长度的整数。但是对于64bit都无法表示的文件怎么办呢？ 二进制协议始终有扩展性不佳的问题。如果是HTTP协议，不管文件多大，只要这个数字还能表示出来，就没有问题。

我们可以看到， HTTP 返回的除了文件大小，还有文件类型等信息，如果是二进制的协议，大概会如图 \ref{fig:binhttp} 这么设计，其中每个方块代表4个字节。要传输的文件紧挨着HTTP头的后面。

\chatu[width=0.2\textwidth]{binhttp}{二进制的HTTP协议，服务器返回部分}

每个类型的信息都分配了固定的字段。解析起来非常快，似乎够用了

等等，我还要添加个信息，表示这个文件该不该被缓存。如果这个文件是临时文件，我大概不希望浏览器将它缓存，而是每次都到服务器上获取文件内容。
好吧，我再添加个字段。

\chatu[width=0.2\textwidth]{binhttp2}{二进制的HTTP协议，服务器返回部分，版本2}

等等，已经运行了的大批服务器和浏览器怎么办？由于文件紧跟着HTTP头，使用新的二进制HTTP协议的服务器碰上老版本的浏览器，缓存属性不仅是被忽略了，更被当做文件内容的一部分了。难道为了一个小修改需要为新的协议修改所有的浏览器和服务器么？

能无限的扩展并保持极佳的兼容性和人类可读性\footnote{测试问题的时候不需要特殊的工具去解析二进制流的各个字段，只需要打印出文本就可以了。}，是文本协议的最大优点。

\section{懒惰原则}

UNIX认为人的时间比机器的时间要宝贵的多。如果有什么事情是机器能代劳的，就绝对不要浪费人的时间。

这是 \textit{The Art of Unix Programming} 的《附录D》收录的无名师的传记，说明懒惰的人写脚本来解放自己的时间。我想没有比这个故事更能表达懒惰原则的含义了。

\begin{insertnote}
\subsection*{无名师与shell脚本}

无名师曾对来访的程序员说：“Unix传统上认为，一行 shell脚本胜过万行C程序。" 

这个程序员自以为对C极其精通，说：“这不可能。Unix内核正是用C实现的。” 

无名师回道：“确实如此。不过，Unix传统上认为，一行shell脚本胜过万行C程序。" 

程序员颇为沮丧：“但是在C中我们可领会到尊者Ritchie的智慧。我们与操作系统和机器合而为一，可以获取无与伦比的性能。“ 

无名师回道：“诚如你言。不过，Unix传统上认为，一行shell脚本胜过万行C程序。” 

程序员冷笑着想悲愤离去。无名师向学生Nabi颌首示意，Nubi在黑板上写下一行shell脚本，问道：“尊敬的程序员，看看这行管道。用纯C实现，是 不是要一万行C代码？“ 

程序员沉吟念诵。最终他承认如此。 

“你需要多长时间来实现和调试那个C程序?"Nubi问道。 

“很长”，来访程序员承认。“但傻子才会干这个而不去完成更有价值的任务。” 

“那么谁更了解Unix传统？”无名师问道。“是写一万行代码的，还是看到任务的无谓而不去编码的？” 

听到此，程序员眼中一亮。
\end{insertnote}

%这个原则我们将在 \ref{section:pkgmgr} \nameref{section:pkgmgr} 一节中讲解autotools的时候体会到。
无名师告诉我们，能用脚本完成就不编写程序。这是懒惰的原则，这是“宁花机器一分不花程序员一秒”的原则。这条原则对非程序员也适用的，就是\textbf{尽量编写脚本将自己的工作自动化。}

%而彻头彻尾的反UNIX，如果让每个软件都倾向于包含你可能用到的一切功能，试图提供一个“家长式”的程序。

\section{UNIX和非UNIX}
UNIX设计哲学并非只能在UNIX身上找到。只是有的系统只学了个表面。更有反UNIX而行之的。

UNIX管道自发明后，已经被各个操作系统借用，包括DOS和Windows都实现了管道。Windows/DOS依旧借鉴了使用|作为管道的办法。
\par\begin{notice}
事实上DOS是作为运行于8086芯片上的蹩脚UNIX开发的。树形目录结构（第二版开始支持）、命令行操作、管道，这些都带有深刻的UNIX影子。
DOS没学到的（受限于8086的能力也无法实现）：多用户多任务，丰富的IPC接口。
\end{notice}

%但是Windows下管道是如此的难用，以至于

UNIX“一切都是文件” 的思想也在许多的操作系统上实现。只是学的并不好。Windows下硬盘实际对应文件为 \texttt{\textbackslash\textbackslash.\textbackslash{}PhysicalDrive1} 第二块硬盘最后一个数字为2。但是这样的文件是无法在文件管理器里看到的。
而且也不能作为文件名使用一般的程序打开。

UNIX使用轻量级进程，进程的创建开销非常小。鼓励开发多个小程序并使用IPC相互协作。
其他系统也实现了多进程能力，但是进程的创建开销（相对UNIX来说）非常大，因而总倾向于使用（便宜的）进程内通信------将当个的程序做的非常大，使用复杂的API进行内部通信。而这些都是BUG的温床。

UNIX鼓励使用文本作为协议。就连配置文件都是（带格式的）文本文件。而有些系统使用专有的二进制格式数据文件保存配置。一旦配置文件损坏，系统就无法使用简单的文本工具进行修复。甚至因为系统本身无法启动而不能进行修复。

\else
\addtocounter{chapter}{1}
\fi 

\chapter{黑客和Linux}

\begin{quotes}{Linus Benedict Torvalds}
If 386-BSD had been available when I started on Linux, Linux would probably never had happened.
\end{quotes}

\begin{quotes}{Richard Matthew Stallman}
Writing non-free software is not an ethically legitimate activity, so if people who do this run into trouble, that's good! All businesses based on non-free software ought to fail, and the sooner the better.
\end{quotes}

在计算机发展的初期，软件\footnote{那个时候系统软件都是汇编语言开发的，所以也谈不上什么源代码。}被作为硬件的附属品销售。
购买一个机器的硬件也就连带的附送了一系列的软件。谁也不会觉得软件是可以被“保护起来”盈利的\footnote{机器码可以非常容易的反汇编重新生成源汇编文件。既然软件是汇编开发的，那也没有什么秘密可以隐藏。}。相互交换彼此的程序随处可见。
直到有一天，西雅图的一个软件天才\footnote{为了避免反应迟钝的人不知道，这位天才是比尔\textbullet盖茨啦}说，他花了很多时间写好的软件，不希望被人随意拷贝。每一个获得拷贝的人都应该给他付钱。
有了天才的带领，自此商业公司开始独立销售软件而不再作为硬件的附属品。

不过这股软件私有化的浪潮还不曾传播到社会的各个交流。
大学里，实验室里，人们还是保持着将软件自由传播的传统。虽然他们不再能自由的传播商业软件，但是自己编写的各种
稀奇古怪的软件可以自由的交流。在没有internet的年代他们就开始利用Usenet\footnote{%
\textbf{来自维基百科：}Usenet是一种分布式的互联网交流系统，它的发明是在1979年由杜克大学的研究生Tom Truscott与Jim Ellis所设想出来的。Usenet包含众多新闻组，它是新闻组（异于传统，新闻指交流、信息）及其消息的网络集合。 }相互交流了。
%
这样一群天才不满于商业化的软件，时常有人对商业公司的系统进行破解，故得名hacker，中译“黑客”\footnote{%
由于只有身怀绝技的聪明人物才有可能破解商业软件。故而黑客给人一种智商非常高的印象。久而久之智商很高（或者因为专业技能好，看起来很高）的程序员都被称作黑客。所以现在说的黑客一般指的是计算机方面非常厉害的人而不再是搞破解的人。现在再搞破解和破坏的人，就没有黑客头衔了，而是Cracker，骇客。
}。

Linux的诞生就是因为黑客门不满足于被大公司控制的操作系统市场。

\section{从UNIX到Linux}

早期的UNIX鼓励源代码的密切交流。这和黑客们的思想如出一辙。伯克利开发的BSD更是吸引了大量的黑客进行贡献。70年代到80年代早期，黑客门过着自由的日子。


但是一场突如其来的官司让AT\&T从反托拉斯的禁锢中挣开\footnote{参考前言的内容。}。
脱缰的AT\&T将UNIX推入了深渊。

但是就算UNIX被推入了深渊，BSD还是一个避风港。但是这最后一点希望都随着AT\&T对BSD的起诉而消失了。

1983年，AT\&T发布了Unix最新版system V UNIX，这是一个商业化版本，付费才能使用，并且不得传播源码。
为了减少纠纷，伯克利分校规定，BSD本身依然保持免费，但是只能提供给持有AT\&T源码许可的公司。
不过，与此同时，伯克利的师生也着手开始将AT\&T的代码从BSD中逐渐去除，以便发布一个不受AT\&T控制的版本。
这一举动惹恼了AT\&T。AT\&T起诉BSD侵犯了Unix的版权。
这场诉讼对BSD打击极大，所有的开发活动都被迫停止，用户人心惶惶，担心自己也遭到AT\&T的追究，因此BSD的使用急剧减少。

这场战争给BSD带来了毁灭性的打击。


黑客们突然间没有了玩具。

最后在1994年，双方达到和解，BSD才恢复开发。但是历史的已经不会再给BSD机会了。因为在BSD官司缠身的时候，上帝把爱心悄悄给了Linux。

1991年，芬兰的一个大学生在网上公布了Linux，立即吸引了大量黑客的目光。从此BSD神码的都成了浮云。黑客门团结在天才周围缔造了Linux的神话。
不用说，这位天才的大学生名字就叫做 Linus Benedict Torvalds。

\section{Linux诞生}

UNIX自诞生以来就是学术界研究的对象。UNIX被移植到大量的机器上，唯独没有染指PC市场。UNIX的人轻视了这个以8086为核心构造的个人电脑。
也难怪，UNIX的多任务能力需要CPU能提供一些保护机制，防止未经授权的访问。而这正是8086所缺乏的。

但是当80386发布的时候，情况有了一些变化。80386具备了“真正的”CPU的特性，包括特权级别区分和分页内存支持\footnote{关于分页技术，请参考 \faqref{FAQ:Paging}。}。 有了这些支持，PC终于具备了运行UNIX的条件。

\subsection{被UNIX忽略的PC}
%那么我们还是字为何PC
计算机最初是作为美国军方的项目研究制造的。美军向来有向全世界开放自己的研究的传统\footnote{Internet也源自美国军方的研究。还有现在每天都在用的GPS卫星定位哦~~。}，所以计算机得以被商业化。
一家造卡片统计机\footnote{恩，先进的自动化的...算盘}的公司因卡片机的大量销售而积累雄厚的财力和强大的销售服务能力，决定进入未知的电子计算机世界。
并于1944年，开发出自动顺序控制计算机，也被称为Mark  I。

\chatu{ibm_mark1}{Mark I计算机}

随后在1951年它决定开始开发商用电脑，聘请冯·诺依曼担任公司的科学顾问，1952年12月研制出第一台存储程序计算机
\footnote{也就是第一台冯·诺依曼结构的电脑。别和第一台（插拔线路编程的）ENIAC混淆哦！}，也是通常意义上的电脑。

没错，这家传奇公司就是IBM。
它开发的第一台电脑就是IBM 701。这是IT历史上一个重要的里程碑。从此计算机进入了程序控制时代，也就是冯·诺依曼结构。

IBM在商用电脑上获得了巨大的成功。那个时候，IBM就是计算机的代名词，有人亲切的称呼IBM为“蓝色巨人”。IBM生产的都是商用电脑，而且都是巨无霸（参考上面的图片），没有一款是适合家庭使用的。
微处理器\footnote{将计算机核心运算和控制部件做到一个芯片上。由于核心非常小，故而得名微型处理器。典型的微型处理器就是Intel为PC生产的80x86系列CPU。}的诞生创造了Apple这样神话。IBM迫切需要一台使用“微处理器”的家用电脑以和Apple这样的新兴公司对抗。

时间紧迫，IBM来不及自己开发操作系统，于是将操作系统外包给了一个西雅图的软件公司\footnote{微软啦！}。CPU也一改自己生产的传统，采购了Intel公司的8088微处理器\footnote{和8086使用相同的指令集，只是在关键的性能指标上比8086次点，因此也有好处：便宜。}，最终的结果就是组合出了一台家用微型电脑。俗称PC（Persional Computer 个人电脑）。

PC就是装配着Intel生产的8088/8086处理器，运行着微软的MS-DOS操作系统的微型计算机。

%\subsection{GNU HURD 内核}
但是8086处理器实在赢弱，缺乏实现多任务所必须的功能。所以DOS是个单任务单用户的操作系统。

\begin{insertnote}
\subsection*{小知识：UNIX需要什么样的CPU}
至少要具备2种功能：特权级别和非特权级别划分、支持内存分页（也就是有MMU，Memory Management Unit内存管理单元）。

UNIX将操作系统内核放到特权级别执行，应用程序则使用非特权级别。当处理器处于非特权级别的时候，对内存和IO的访问是严格受限的：访问不允许的IO设备会产生一个异常，该异常被操作系统捕获\footnote{所谓异常捕获，就是处理器发生异常的时候自动跳转到内核提供的异常处理代码。}，然后中止试图越权访问的程序；访问不允许访问的内存——如内核使用的内存——会导致内存访问异常，操作系统捕获这个异常，中止试图越权访问的程序。内核自身运行于特权模式。
有了CPU的这种级别划分，UNIX能保护操作系统不受异常程序的影响，提高稳定性和安全性。

UNIX支持多进程运行。每个进程\footnote{参考 \faqref{FAQ:Process}}之间的内存必须进行隔离。也就是说，一个进程只能访问自己的内存，无法访问其他进程的内存。
事实上这个访问不是通过“权限检查”（试图访问其他进程的内存导致进程被中止）实现的，而是通过“完全的隔离”实现的。
完全的隔离，也就是任意进程拥有完整的内存地址空间。并不存在内存地址被划分成N个区域，每个进程只能访问属于自己的$ \frac{1}{N}$的区域。而是每个进程拥有整个的地址。
从 0 开始到 $2^{32}-1$ （对应于32位的系统，如果是64位的则是 $2^{64} - 1$），访问的通通都是属于该进程自身的内存。
不同进程可以使用相同的内存地址访问属于自己的不同的数据。这种内存隔离被称呼为“虚拟内存”。虚拟内存是支持多进程运行必备的功能。

虚拟内存必须依靠CPU的MMU提供的分页功能才能实现。
\end{insertnote}

由于8086/8088处理器并不拥有MMU，也没有实现特权级别划分，所以根本不具备运行UNIX的条件。

就这样PC市场被UNIX忽略，直到80386的出现。80386是Intel生产的第一个具备运行UNIX自个的CPU。
但是就算在功能上具备了运行UNIX的资质，80386从UNIX的观点来说，实在是太慢了。所以谁也没打算将UNIX移植到使用80386处理器的PC上。

PC机就这样继续运行着DOS，以及后来的Windows。

\subsection{386BSD、 Minix催生Linux}

要说PC真的没有一个UNIX可以运行那是假的。商业UNIX抛弃了PC，但是黑客们却像发现了新大陆一样为PC着迷。于是，各种黑客版本的UNIX被移植到80386芯片上。
其中最著名的莫过于386BSD（将BSD移植到80386处理器的项目）和Andrew S. Tanenbaum教授以教学为目的开发的Minix。

Linus在大学里用过UNIX就为它着迷。于是决定为自己买一台计算机运行UNIX，这样就不要跑去学校的机房了。
由于PC性能不佳，跑着UNIX的工作站又太贵了。Linus决定购买一台PC。时值BSD正好处于和AT\&T的长期诉讼中。法律风险是大大的。于是Linus装起了Minix。
至于其他版本的UNIX，就没入Linus的法眼，呵呵。

Minix是Andrew S. Tanenbaum教授以教学为目的开发的UNIX，但是并不是一个“自由”的软件。而且Linus对于Minix有相当多的吐槽。不仅仅是Linus，许多人都为他们的吐槽写了补丁。
但是Andrew以“只用于教学目的的系统不想变的太复杂”为由而拒绝添加这些改进。一个拒绝改进的开发者+虽然宽松但是不自由的授权\footnote{
购买Andrew的《操作系统设计与实现》这本书就能获得Minix代码并授权个人使用。仅限于非商业用途。}\textrightarrow ~Linus~决定\textbf{为自己}开发一个UNIX兼容的操作系统。


\subsection{以GPL开源}

Linus说自己在Linux上做的最正确的决定就是用GPL保护了Linux。1911年发布的第一版内核并不是使用GPL协议发布的。
他使用了自定的一个协议，内容大概就是：{\it Linux自由免费，随便用。如果你对Linux做了修改，麻烦你把修改后的代码给我一份，谢谢。}

在黑客界，有一个古老的传统：\emph{尽量去发现工具，实在没有就自己写一个。} 这条传统似乎对软件授权协议也有效。
%
Linus大概觉得自己跳过了第一步的发现，于是他决定重新开始依照老传统做一次——这次他发现了GPL，于是他不再发明，将Linux置于GPL的保护之下。

Linus为第一版Linux编写的协议和GPL想要表达的意思差不多，但是起草过程中有律师帮忙的GPL更加的严谨。

从头开发的Linux源码和GPL授权为Linux的发展构建了坚实的后盾。即没有侵权的后顾之忧也没有不怕被侵权后缺乏法律的保护。Linux就这样替代了386BSD，Minix等等等等
成为了PC上最受欢迎的类UNIX系统。同时由于Linux使用GPL授权，填补了GNU缺失的内核 GNU Hurd。

全世界的黑客开始行动起来，推动着Linux快速的发展。Linux成了分布式开发的典范。

\chapter{初识Linux}

\begin{quotes}[0.63]{Linus~Torvalds}
To kind of explain what Linux is, you have to explain what an operating system is. And the thing about an operating system is that you're never ever supposed to see it. Because nobody really uses an operating system; people use programs on their computer. And the only mission in life of an operating system is to help those programs run. So an operating system never does anything on its own; it's only waiting for the programs to ask for certain resources, or ask for a certain file on the disk, or ask to connect to the outside world. And then the operating system steps in and tries to make it easy for people to write programs.
\end{quotes}

%\subsection{使用Linux内核的操作系统——发行版}

\section{Linux是操作系统内核}

TODO：Linux 主要发行版介绍，何为 Linux 发行版 Linux 发行版的作用

\section{发行版的意义}
\section{包管理}
\section{和软件仓库}


\section{包管理和软件仓库}

\subsection*{rpm 和 deb}
\subsection*{Portage}


\chapter{快速安装 Linux}
Linux介绍了那么多，是时候将它安装到我们的电脑上了。要安装一个Linux系统，我们总是选择一个发行版来安装。那么到底选择什么样的发行版呢? 理想的发行版通常是这样的一个发行版：它非常第容易安装，非常的容易适用；软件仓库里的软件非常丰富，我可以用包管理安装一切软件而不需要自己手动安装；对于一些软件我希望能容易的进行定制；快，非常快。在对各个发行版做里权衡后，我选择Gentoo作为本书中介绍Linux知识所使用的平台。


\begin{notice}
 注意：Gentoo只是我们选择的众多发行版中的一个，我并不打算把本书变成一个Gentoo入门教程。如果这样，读者还不如选择阅读Gentoo官网上的手册。我尽量将所涉及到第知识通用化。如果这样的知识点只适用于Gentoo, 我会进行说明并辅之以其他发行版的等位操作\footnote{不同发行版达到同样的目的采取的不同操作。}。
\end{notice}

本书并不打算引起发行版之争\footnote{指不同的人使用不同的发行版，并常常在网上争论自己使用的发行版要比别人的优秀}，
但是又必须解释为什么使用了Gentoo这个非常罕见的发行版而不是同其他书本那样采用最流行的发行版，比如Ubuntu或者Red Hat系。

最直接的原因也许就是作者本身使用的就是Gentoo，以自己最熟悉的发行版写作，可以避免不熟悉系统进行错误的讲解\footnote{这确实是一个理由，但是Gentoo必定有比其他的发行版更值得作为本书的平台。}。
当然，如果仅仅这一个理由未免太牵强。
\it
Gentoo可以让读者将精力放到学习Linux本身，而不需要关注发行版的细节，这就是我采用Gentoo的原因。
\normalfont
另一个能做到这种效果的发行版为LFS\footnote{严格来说，LFS 算不上发行版，只是一个教你如何手工编译出一个能用的系统的手册。}（Linux From Scratch）。
但是LFS太繁琐，恐怕还没有学习到Linux的知识就已经中途放弃了。


\section{安装准备}

当然是准备一台电脑啦！Gentoo虽然号称为本机编译优化，但是编译本身非常耗费系统资源，所以一台主流配置的电脑是必不可少的！
\footnote{Gentoo支持的硬件架构非常多，不是一定要用个人电脑的。
但是作为初学者，我们就用最常见的个人电脑，不折腾各种稀奇古怪的电脑，最重要的是，Gentoo对x86体系也是支持的最好的（这不废话么，用户量摆那里么）。}


除了一台电脑，还要准备个Fedora或者Ubuntu的LiveCD\footnote{参考附录A. FAQ 的“什么是LiveCD”。
}，这里不推荐Gentoo自己的LiveDVD\footnote{同LiveCD。}，因为前者更容易获取和使用。

当然，还有网络。Gentoo虽然可以脱离网络进行安装，但是需要将所可能需要用到的源码提前下载，这是非常麻烦的。所以最好准备一个安装的时候可以访问的网络，也就是说LiveCD 环境里必须能使用网络。Fedora 或者Ubuntu那样的LiveCD 通常提供了完整的桌面环境，使用这些桌面环境的网络设置，让LiveCD 能访问互联网就可以了。你可能需要使用 NetworkManager 提供的图形客户端进行配置，通常是在桌面右上角的系统通知区域。要使用 NetworkManager 配置网络，请参考第6章的内容。

嗯嗯，不过就算是x86还分为x86\_32（ 没有特殊指明的情况下x86特指x86\_32） 和x86\_64\footnote{因为历史原因，AMD率先推出支持64位运算的x86CPU，所以Gentoo管x86\_64叫amd64，这和x86又被称呼为 i386 是一样的原因。}呢。
由于x86\_64是在x86的基础上发展起来的，所以x86\_64向下兼容x86。 也就是说，64位的CPU可以运行32位的软件。但是反过来就不行了。

既然存在两种x86，选择起来就费脑筋了。怎么知道自己的CPU是x86\_64的还是x86\_32 的呢？答案是，现在生产的CPU都是x86\_64了。当然如果你想准确的知道CPU到底支持不支持64位指令，在Window下你可以选择执行CPU-Z 软件获得CPU信息，如果是Linux就简单的多，在终端下执行lscpu命令就可以了。

（1）  Linux下使用lscpu：

在终端下执行lscpu，该命令的第一行输出就是CPU的体系结构。在我的电脑上，其输出为：

\begin{code}%
Architecture:        x86\_64\\
CPUop-mode(s):        32-bit, 64-bit\\
Byte Order:            Little Endian\\
CPU(s):                8\\
On-lineCPU(s) list:   0-7\\
Thread(s) per core:    2\\
Core(s) per socket:    4\\
Socket(s):             1\\
Vendor ID:             GenuineIntel\\
CPUfamily:            6\\
Model:                 42\\
Stepping:              7\\
CPUMHz:               1600.000\footnote{如果CPU支持动态降频，这里显示的是当前频率而不是CPU的标称频率。}\\
BogoMIPS:              6420.96\\
Virtualization:        VT-x\\
L1d cache:             32K\\
L1i cache:             32K\\
L2 cache:              256K\\
L3 cache:              8192K\\
\end{code}

第一行的Architecture:        x86\_64就表明此CPU的结构是x86\_64。


\vskip 1em
（2）  Windows 下使用CPU-Z：

\chatu{cpuz}{cpuz截图}
如图\thefigure所示。我红圈标注的EM64T表示CPU支持64位指令。对于AMD系的CPU，您大可放心，全部支持64位指令。


\begin{insertnote}
\subsection*{小插曲  64位 vs 32位}

如果是x86\_64的CPU，应该选择amd64\footnote{AMD64 就是x86\_64。}
的Gentoo还是x86的呢？既然现在的CPU都支持 64位指令的，为何还有那么多人使用 32 位的系统？

x86\_64虽然能执行x86代码，但并不是没有代价的。Windows下大部分现有的程序都是32位的，如果安装的64位Windows，32位软件并不能100\%兼容。而且所有的64位软件同时有32位的版本，所以Windows用户非但没有迫切的升级到64位系统的意愿，放而有为了兼容性考虑不得不安装32位系统的处境。

64位系统如果不执行64的程序，那性能还不如直接使用32位的系统和32位的程序。所以在 Windows 统治的世界里，64 位CPU的地位非常尴尬。32位CPU通过PAE\footnote{Physical Address Extension，物理地址扩展。允许32位CPU使用36位地址从而突破4G内存限制。}同样能获得 >4G 内存的支持（非服务器版Windows人为阉割32位版本的PAE支持，故而32位桌面版Windows不能支持超过4G内存。）。

大概也就是因为这个原因，所以不明真相的群众就以为64位除了支持大内存外一无是处。继续固守32位。64Bit vs 32bit的战斗因此打响。

x86\_64对x86的提升不仅仅是更大的内存，同时增加的有

(1)  更大的寄存器 64位CPU自然拥有64位宽度的寄存器。

(2)  更多的寄存器x86\_64相比x86增加了8个通用寄存器,增加8个MMX寄存器。

(3)  默认支持SSE\footnote{SSE指令是Intel为x86添加的扩展，加快了浮点运算的速度。}指令集。

(4)  更有效的指令编码。

这些都无形中增加了x86\_64指令编译的软件的运行速度。所以，如果你的CPU支持x86\_64，请必定选择x86\_64。
Linux上没有Windows上遇到的问题。
Gentoo下几乎所有软件都是本机编译的，自然不存在兼容性问题。
而个别非开源软件也多数拥有 64位版本，更不需要操心兼容性问题。
\end{insertnote}



好了，将LiveCD塞入光驱，开始一步步安装属于你的Gentoo吧。有关如何使用LiveCD请参考\faqref{FAQ:UseLiveCD}

好，我们列一个清单：


\begin{itemize}
\item[ \checked] 一台电脑
\item[ \checked] 一张LiveCD
\item[ \checked] Internet
\item[ \checked] 本书或者Gentoo 安装手册。
\item[ \checked] 耐心，很多很多的耐心
\end{itemize}

你准备好了么？

\section{分区规划}
任何系统都是安装到硬盘使用的
\footnote{好吧，先忽略掉LiveCD 和  WinPE 这类不需要安装的系统，咱讨论的是一般用途的桌面操作系统。}，
安装到硬盘之前，必须先划好家。在介绍如何分区前，首先得知道什么是分区，然后参考
FHS\footnote{FHS 是文件系统目录结构的一个标准。规定了根分区下各个子目录的名称和用途。}
指示结合自己的实际情况规划好分区。

\subsection{分区基础知识}
最初，计算机使用软盘，软盘是没有分区的。后来蓝色巨人IBM发明了硬盘。硬盘的容量一下子比软盘大出好多倍。当时的MS-DOS所使用的FAT12文件系统无法管理那么大的硬盘。于是西雅图的巨人和蓝色巨人想出来把硬盘划成逻辑上的几个区域，每个区域大小都在MS-DOS能管理的范围之内——这样的逻辑区域就是分区。人们发现将硬盘划分为逻辑上的几个区域后，更容易组织硬盘上的数据了；而且一个文件系统错误只会影响到一个分区的数据，其他分区不受影响，数据的安全性也得到的提升。因而后来的DOS虽然将FAT12进化到了FAT16，能管理当时的大容量硬盘了，但是分区这个功能却保留了下来\footnote{虽然是MS-DOS的发明，但是Linux可不会拒绝这样的发明。Linux还支持 BSD 发明的分区格式，总之，Linux决定支持越多的分区表格式越好，这极大的方便了用户，不是么？}。

既然用户划了分区，操作系统总得知道用户到底怎么划分的，描述分区的数据被称作分区表。既有分区表，必须有个地方存储，也必须知道到哪里去读取分区表。MS-DOS的把分区表和引导程序放入硬盘的第一个扇区\footnote{扇区是硬盘最小寻址单位。参考附录A. FAQ里的条目“什么是扇区”。}。

硬盘的第一个扇区又被称呼为MBR\footnote{MBR是硬盘的第一个扇区，具体解释请参考附录A. FAQ 里的条目“什么是 MBR”。}。
MBR既要存储引导程序，又要储存分区表，是个寸土寸金的地方，分区表大小受限，只有4个表项。也就是说，一个 MBR分区表最多只能有4个分区。要是只有四个，似乎并不够用。

MS-DOS将其中一个分区作为扩展分区，然后再扩展分区里再建分区表。扩展分区里面的分区就是逻辑分区，MBR 上的分区表就是主分区表。主分区表里划分的分区自然就是主分区了。所以一个MBR格式分区表最多允许3个主分区和1个扩展分区，或者4个主分区。扩展分区里再创建逻辑分区，没有数量限制。
\chatu{logicalpart}{链式逻辑分区表}

\begin{notice}
注意：扩展分区里的逻辑分区并不是表格形式存储的，而是“链式”存储。如图\thefigure所示。每个逻辑分区包含查找下一个逻辑分区的信息。因而一个逻辑分区的破坏有可能造成链式效应，将所有的逻辑分区全部摧毁。
\end{notice}

由于近来UEFI\footnote{参考附录A “UEFI 和 BIOS”。}  的兴起，
UEFI指定的分区表格式GPT也流行开来了。
和MBR不同的是UEFI不需要专门的引导扇区，引导程序由UEFI直接从文件系统上加载。所以GPT只是分区表，不需要和引导程序共存。

MBR 的分区表只有64个字节大小，只能包含最多4个分区的信息。
GPT包含128个分区表项，最多允许一个硬盘划分成128个分区，足够了。
能表示更多的分区并不是GPT唯一的优点，MBR分区表只能管理2.2TB
\footnote{MBR使用4个字节表示分区起始位置的偏移量。
偏移量以扇区为单位，一个扇区为512字节，那么MBR分区表能管理的最大硬盘大小为$2^{31}*512Byte = 2TB$。
又因为硬盘厂家以1000为进制而不是1024，故而大约为2.2TB。}%
以下大小的硬盘，GPT却可以管理容量超过9ZB\footnote{1ZB=1024PB 1PB=1024TB. 硬盘厂家的计算是 1ZB=1000PB 1PB=1000TB。}的硬盘。
在大容量硬盘越来越普遍的今天，MBR显得越来越力不从心。GPT正好接替MBR成为今后PC硬盘的主流分区表格式。

如果你电脑的固件是UEFI，不管硬盘实际是否大于2TB，建议最好使用GPT分区表。

\subsection{为磁盘分区}

\chatu{gparted-gui}{gparted界面}

\FloatBarrier

\normalfont

有很多种工具可以对硬盘进行分区。今天的主角是 gparted , 这个是一个非常简单易用的分区管理软件，带图形界面的哦~~

要打开gparted，在终端执行sudo gparted。如果 Fedora 的LiveCD提示没有这个命令，则用命令sudo yum install gparted或者apt-get install gparted进行安装。gparted界面如图\thefigure所示。

在gparted里所做的一切操作都不是立即执行的，只有应用操作后才被一次性执行。所以未应用前撤销操作是来得及的。
下面我把为硬盘分区所需要的几个主要操作图解一下。请读者根据此处的图示熟悉操作流程，并在学习下一个小节后自行分区。

\chatu{gparted-apply}{gparted应用操作}
\noindent要应用操作：\\\indent
1. 单击“编辑”|“应用全部操作”菜单。如图\thefigure所示。\\\indent
2. 在弹出的警告对话框中选择“是”。

\chatu{gparted-newpart-boot}{gparted新建分区}

\noindent建立一个分区：\\\indent
1. 选择磁盘的空白区域，右击后在菜单里选择“新建”。\\\indent
2. 在弹出的对话框里指定一个大小、分区类型和文件系统，如图\thefigure所示。\\\indent
\hskip 4em %
\begin{notice} %
  注意：分区表为GPT格式的只能选择主分区。
\end{notice} \\\indent
3. 单击“添加”按钮，新建的分区就编排在了任务列队里了\\\indent

\chatu{gparted-delpart}{gparted删除分区}
\noindent删除一个分区:\\\indent
1. 右击要删除的分区在菜单里选择“删除”。如图\thefigure所示。 %\\\indent

\chatu{gparted-newpt1.png}{新建分区表}

\noindent新建一个分区表：\\\indent
1. 单击“设备”｜“创建分区表”菜单，如图\thefigure所示。

\chatu{gparted-newpt2.png}{新建分区表}

\indent2.在打开的对话框里点“高级”前面的箭头，选择一个分区表类型。默认为MS-DOS分区表，也就是MBR分区表。
我需要在这里我选择gpt分区表。如图\thefigure所示。

\FloatBarrier

好了，gparted的基本操作就完成了，接下来学习一下文件系统的结构后按照自己的意愿为自己的系统分区吧。

\subsection{文件系统结构标准（FHS）}

Linux继承自UNIX树形目录。每个目录各司其职，并被FHS（Filesystem Hierarchy Standard）标准化。表3.1所列是一个不怎么完整的FHS标准，但是对于读者决定如何为他自己的电脑分区已经足够了。

\begin{longtable}{|l|p{0.7\textwidth}|}
\caption{FHS标准参考}\\
\hline
/ & 根目录 \\\hline
/bin & 基本系统程序。如 ls cat grep \\\hline
/sbin & 只供管理员使用的基本系统程序 \\\hline
/boot & 内核和引导程序，通常为独立分区\\\hline
/dev & 设备文件目录 \\\hline
/etc & 存放系统配置文件。不允许存放为独立分区\\\hline
/home & 存放非root用户的家目录。 
比如foo用户的家目录就是/home/foo
一般用来存在个人文件和个人设置。
通常强烈建议独立为一个分区。\\\hline
/lib & 系统基本库，被 /bin /sbin 里的程序依赖的库 \\\hline
/media & 可移动媒体的挂载目录。比如cdrom和u盘。被/run/media代替\\\hline
/mnt & 临时挂载目录\\\hline
/opt & 通常用于安装非开源软件，如Adobe Reader \\\hline
/proc & 虚拟的文件系统，用来获得内核和运行中的程序的信息\\\hline
/root & root用户的家目录\\\hline
/srv & 做服务器使用的时候 ，用来存储服务的数据。如一个git服务器通常将仓库存放于/srv/git\\\hline
/tmp & 临时目录。其中文件关机丢失。建议挂载为tmpfs虚拟文件系统，对于tmpfs的介绍可以参考第7章。 \\\hline
/usr & 非基本系统程序根目录，下面的结构和 / 差不多 \\\hline
/usr/bin & 同 /bin。只是非基本程序。如xeye  elinks . \\\hline
/usr/sbin & 同 /sbin。只是非基本程序。如system-config-firewall \\\hline
/usr/include & 头文件目录。\\\hline
/usr/lib & 同 /lib。只是其包含的库是不被/sbin、/bin使用的。\\\hline
/usr/share &  保存架构无关的共享数据。如图标\\\hline
/var & 系统运行中会不停变化的文件。比如各种log，包管理器的数据库，等等。通常在服务器上是独立分区，个人电脑可以不独立划分。\\\hline
/var/cache & 缓存。该目录下的文件可以安全删除，要求使用它的程序必须能重建它。比如fontconfig的缓存。\\\hline
/run 和 /var/run & 临时目录。包含本次系统运行时的信息。\\\hline
/var/log & 存放各种日志文件。\\\hline
/var/tmp & 临时文件，可能重启后还在。注意和/tmp的区别。\\\hline
\end{longtable}


知道了系统各个目录的作用，就好为目录进行分区了。通常我们使用4个分区 /boot、/、/home 和一个交换分区。

\begin{notice}
注意：我在使用“/home分区”这个术语的时候，我指的是为“/home”分一个单独的分区，并且将这个分区挂载到“/home”目录下。我通常不使用“挂载到/home目录的分区”而直接简述为“/home分区”。同样的，“/分区”（或者根分区）指的就是挂载到根目录的分区。
\end{notice}

/boot独立为一个分区的好处是防止根文件系统的错误蔓延导致/boot/*文件破坏无法正常启动。另一个好处是 /分区可以使用引导程序（如GRUB）不支持的文件系统格式。例如GRUB不支持XFS文件系统格式，但是你可以让/boot分区使用ext2格式从而使 /分区可以使用XFS格式。如果/boot没有独立出一个分区。那 /分区就只能使用引导程序所支持的文件系统了。

/home独立一个分区的好处就是将用户数据彻底和系统数据分离。由于/home目录通常频繁读取写入，和 /分区隔离可以避免 /分区被碎片化。

提示：其他目录，比如“/var”，可以独立一个分区也可以不独立。就看你系统的用处了。作为桌面用途，完全没必要。如果是服务器，“/var/log”可能会很大。邮件服务器的“/var/mail”就可能会非常大，所以非常有必要使用独立分区。

\subsection{文件系统选择}


Linux支持的文件系统非常多，有自家系列 ext2/ext3/ext4、甲骨文公司开发的btrfs，来自硅谷图像的XFS等等等等。

不过选一个文件系统没那么纠结的。对于根分区\footnote{根目录所在分区}，既然大部分发行版默认使用ext4，那就ext4吧。

为了引导程序的兼容性，/boot分区仍旧使用ext2或者ext3。

对于/tmp这样的目录，强烈建议使用tmpfs\footnote{参考第7章“文件系统”。}。tmpfs是一种内存文件系统，所有tmpfs上创建的文件都是在内存里的，关机后就会消失。非常适合用作/tmp的文件系统。既能加快系统速度，又能保护硬盘。

对于/home目录，如果是个BT爱好者，会有很多大文件，我推荐使用XFS文件系统。其他情况下还是ext4各方面性能比较平衡。

如果希望为自己选择最佳的分区格式，请参考第7章“文件系统”的内容。

\section{开始安装}

废话了那么多，介绍了一些安装前需要知道的基础知识，现在终于要进入正题了，安装Gentoo。
Gentoo的安装非常简单，首先是分区。这个前面已经介绍过了。我现在假定读者已经在LiveCD 环境中完成分区了。下面开始真正的安装。


\subsection{Shell 快速入门}

由于安装Gentoo需要执行shell命令，所以我需要讲解一下最基本的shell命令。

机器只能接受二进制指令，可是人需要工作在更高一层。人希望直接提供文本形式的指令由机器执行。shell 是操作系统提供给用户的解释器。shell把来自用户的命令转化为机器指令并执行。

命令虽然是文本，但也不是自由格式的。计算机是死的，自然不能接受人类的语言，只能是比较固定的几个格式。

一条shell命令的最常用格式为:

\begin{code}
命令 [选项] [目标]
\end{code}

选项是可选的，可以有一个或多个；多个选项间用空格隔开；有些命令没有选项；目标通常是一个文件的文件名，或者目录名。

请看一个命令做为例子： 

\begin{code}
root@gentoo \textasciitilde \# ls -l --all /  
\end{code}
在这里：

\begin{itemize}
\item -l 是一个短选项。短选项 “-”开头，后面跟单个字母。在这里 -l 表示以“长格式”列出。所谓“长格式”就是相对与一般格式，除了要列出文件名，好要列出文件所有者，权限模式，文件大小
\item -{}-all 是一个长选项。长选项使用两个“-” 开头，后面跟一个单词。在这里 -{}-all 表示列出所有 文件，包括隐藏文件。
\item / 在这里为操作对象参数，在这里是一个目录。表示要列出 / 下的文件。
\item ls 的所有参数都是可选的，如果不使用任何参数，ls 列出当前目录下的文件。


\end{itemize}

shell命令分两种，一种是内部命令，另一直是外部命令。内部命令是由shell直接执行的\footnote{最常用的内部命令是cd。}。
外部命令则是独立的程序。比如ls就是独立程序，它的路径为/bin/ls。shell使用一个PATH环境变量\footnote{参考附录A FAQ“什么是环境变量”。}来查找外部命令。

现在我们只需要知道shell最常用的命令形式就可以了。在以后的学习过程中，我会像大家介绍 shell 的其他命令格式，包括各种循环和函数使用。现在还要一个比知道更多的shell命令更重要的知识需要知道：
\begin{center}\em
shell字符串扩展
\end{center}

还是 ls 这个命令作为例子，假设当前目录下有“a.txt”“ab.txt”两个文件。执行下面的命令：

\begin{code}
root@gentoo \textasciitilde \# ls *.txt
\end{code}

结果就是列出 a.txt ab.txt 两个文件。

在这里：
\begin{itemize}
\item *.txt 会被扩展成 a.txt ab.txt
\item a.txt 和 ab.txt 被作为2个参数传递给ls。
\item “*”是通配符，表示匹配任意多个字符；“?”表示匹配任意的一个字符。
\end{itemize}

比如你执行的是下面这个指令：

\begin{code}
root@gentoo \textasciitilde \# ls ?.txt
\end{code}

结果就只显示a.txt文件。要取消shell字符串扩展，加引号即可，如下面这条命令：

\begin{code}
root@gentoo \textasciitilde \# ls "*.txt"
\end{code}

执行结果为：

\begin{code}
ls: 无法访问*.txt: 没有那个文件或目录
\end{code}


我们知道，ls *.txt 相当于 ls a.txt ab.txt，可是如果执行下面的命令

\begin{code}
root@gentoo \textasciitilde \# ls "a.txt ab.txt" 

ls: 无法访问a.txt ab.txt: 没有那个文件或目录 

root@gentoo \textasciitilde \# 
\end{code}

shell将“a.txt ab.txt”整体（不包含引号）作为了一个参数传递给ls了。

要理解shell的这种行为，首先要明白shell是如何传递参数给程序的。


\subsection{init 系统选择}

\subsubsection{sysvinit+openrc}
\subsubsection{SystemD}

\subsection{内核编译}

\subsubsection{新手使用 genkernel}

\subsection{引导管理器}
\subsubsection{UEFI}

\subsection{安装基础系统}
\subsection{系统初步配置—重启前配置}

\subsubsection{nano 编辑器快速入门}

\chapter{Linux 探索}
\section{Shell 基础和日常应用}

\section{文件管理}	
\section{文件编辑器}
\subsection{编辑器之神 vim	}
\subsection{神之编辑器 emacs	}
\subsection{够用就好 gedit	}
\section{系统管理基础}

\subsection{进程管理	}
\subsection{系统服务管理	}

\section{ shell 进阶应用}

\chapter{软件管理}
\section{ 软件的分发	}
\section{软件的编译	}
\section{发行版的包管理}\label{section:pkgmgr}
	
\chapter{图形界面}
对命令行是不是有点枯燥了？是时候来电新鲜的东西了 只所以现在才讲，因为图形实在是太复杂了。

\section{窗口和图形}

\section{X 的历史}

TODO： X 是一个窗口系统，那么X为什么采用了 C/S 结构呢？
\subsection{ 窗口管理}
\subsection{ 渲染库和静态链接}
\section{桌面环境}
\subsection{GNOME}
GNOME 是事实上的标准。各个主流发行版的默认桌面环境。
\subsection{ KDE XFCE 和其他}

\section{  办公和打印	}
\subsection{  LibreOffice	}
\subsection{使用CUPS 打印文档	}
\section{ 声音的那些事	}
\subsection{  声音的数字编码	}
\subsection{  逐渐被遗忘的OSS	声音架构}
\subsection{ ALSA 高级Linux声音架构}
\subsection{  PulseAudio 优秀的声音服务器	}
\subsection{  播放器	}
\section{  让游戏放松自己	}
\subsection{  OpenGL 和 Mesa	}
\subsection{  硬件加速	}
\section{  Wayland	}
\subsection{ X 的错误	}
\subsection{  直接渲染}

\chapter{网络}
\section{  网络基础知识}
\subsection{  通信基础}
\subsection{  IP地址和TCP/IP协议}
\subsection{  路由表和路由协议}
\subsection{ 以太网和WIFI}
\subsection*{  小插曲：永恒的以太网}
\section{  网络配置}
\subsection{  图形环境下的工具	}
\subsection{ 命令行配置工具}
\subsection{ 网络自动配置和DHCP}
\section{  网络攻击和防火墙	}
\subsection{ 网络攻击类型和检测}
\subsection{  iptables防火墙}

\chapter{文件系统}
\section{  磁盘上的文件系统	}
\subsection{  块设备和磁盘}
\subsection{ 超级块和inode}
\subsection{  典型磁盘文件系统举例}
\section{  虚拟文件系统	}
\subsection{  /proc内核信息窗口}
\subsection{  /dev设备文件系统和udev}
\subsection{  虚拟内存盘tmpfs}
\section{  网络文件系统	}
\subsection{  NFS	}
\subsection{  Windows 网络邻居 CIFS	}
\section{  其他的文件系统	}
\subsection{  LiveCD的最爱——压缩文件系统squashfs}
\subsection{  为Flash芯片设计的文件系统}

\chapter{架设服务器}
\section{  搭建 HTTP 服务器}
\subsection{  apache 用的最多的服务器}
\subsection{ nginx 轻量级服务器}
\subsection{ lighttpd 超轻量级服务器}
\subsection{ squid 加速代理	}
\section{  数据库	}
\subsection{  最流行的开源数据库 MySQL}
\subsection{  最优秀的开源数据库 PostgreSQL 	}
\subsection{  商业霸主 OracleDB	}
\section{  加速 DNS ，在本机搭建 DNS	}
\section{  共享打印机	}
\subsection{ CUPS 打印服务	}
\subsection{ Samba 打印机共享	}


\chapter{压榨机器的性能}
\section{  最廉价的优化 – 编译优化	}
\subsection{  编译期优化	}
\subsection{  链接时优化（LTO）	}
\section{  有目的的优化 – 查找性能瓶颈}
\subsection{ 查找性能瓶颈}
\subsection{  优化瓶颈}
\section{  benchmark – 优化的基准测试}
\subsection{  CPU 单核能力测试}
\subsection{  CPU 多线程性能测试}
\subsection{  文件系统 IO 测试}
\subsection{  网络性能测试}
\subsection{  GPU/OpenGL 性能测试}


\chapter{时光机器-版本控制系统}
\section{   历史和后悔药	}
\section{   中心式版本控制仓库 CVS	}
\section{   CVS 后继 SVN	}
\section{   GIT 划时代的分布式版本控制	}
\subsection{  Bitkeeper	}
\subsection{  Linus 消失的一周	}
\subsection{  版本控制设计为一个文件系统	}
\subsection{ 去中心化	}
\subsection{  GIT典型工作流	}
\subsection{  GIT 简单使用	}


\appendix

\renewcommand{\thesection}{\arabic{section}}

\chapter{FAQ}\label{chap:FAQ}

\section{什么是 LiveCD ?}\label{FAQ:LiveCD}
一种让 Linux 系统脱离硬盘，只需要一张 DVD 或者 CD 运行的技术。使用者只需要插入一张 LiveCD 就可以使用 Linux  系统了。当然，由于 CD-ROM 是只读的，所以对 LiveCD 系统本身所做的更改重启后都会消失。由于 Linux 发行版的膨胀，一张 CD-ROM 无法装下，于是就有了 LiveDVD , 除了容量增加外，和LiveCD 是一致的。此外还有 LiveUSB , 也就是使用 USB 存储设备的 liveCD。用法和 LiveCD 并无二致。

\section{如何使用 LiveCD。}\label{FAQ:UseLiveCD}

\section{UEFI和BIOS是什么？}


\section{什么是API?。}\label{FAQ:API}

\section{Shell 是什么？}
操作系统是由内核和一系列的系统程序组成的。

\section{什么是进程和线程？}\label{FAQ:Process}

\section{什么是内存分页？}\label{FAQ:Paging}

\section{如何不重新登录切换到root帐号？}\label{FAQ:su}

正如 \secref{sec:usersandgroups} 所讲，UNIX支持多账户登录。如果当前使用普通帐号登录，而某些命令需要使用root帐号的时候，就需要用到su这个命令了。
使用su命令的前提是加入wheel组，否则不能使用su切换到root。 “su”切换到root帐号，但是没有切换环境变量和工作目录，亦没有执行/root/.bash\_profile。{}“\hbox{su -}” 能切换到 /root 目录，提供类似直接登录的环境。用户很可能需要的就是带 - 参数的su。

su 还可以切换到其他用户，如“su smith - ”可以切换到smith帐号上。

另外一个临时使用root权限的办法就是在执行的命令前加 sudo 。 sudo 需要用户加入wheel组并在 /etc/sudoers 加入 “\%wheel ALL=(ALL) ALL” 一行。
sudo 通常来说只应该授权给信任的普通帐号 —— 比如管理员自己给自己设定的日常登录帐号（如果是桌面用户，自然就是自己用来登录桌面的帐号） —— 并且为普通用户设定好强壮的密码。

\section{内核参数}\label{FAQ：kernelparamter}


\section{Initramfs是什么}　\label{FAQ:initramfs}

Linux启动过程中内核主动做的最后一件事情就是执行/sbin/init，那么倒数第二件事情必然是挂载根目录。
挂载根目录有2个前提条件：找到根分区和支持根分区所使用的文件系统。
找到根分区，第一点，需要其所在设备的驱动。第二点就是内核需要知道哪个是根分区。这个可以通过内核参数 root= 设定。

Linux下驱动可以编译进内核，也可以编译为模块动态加载。如果恰巧磁盘设备的驱动被编译成了模块，就出现“鸡生蛋和蛋生鸡”的问题了：不访问磁盘就无法加载模块，而不加载模块有没有磁盘的驱动。这对于自己编译内核的童鞋不成问题。只要知道自己的SATA/IDE控制器的型号，将对应的驱动编译进内核就解决了。但是如果是打算制作“通用内核”的发行版呢？

文件系统也是同样的问题，没有文件系统驱动就无法挂载分区访问上面的文件，无法访问文件就无法加载模块。这对于自己编译内核的童鞋不成问题。只要将自己根分区所使用的文件系统编译进内核就解决了。但是如果是打算制作“通用内核”的发行版呢？将所有的文件系统都编译进内核么？

initramfs就是为解决这个“鸡生蛋和蛋生鸡”的问题诞生的。

GRUB加载内核的时候通过initrd 命令加载一个initramfs，该initramfs就会成为内核的根分区\footnote{事实上是Linux内核拥有一个不挂载就存在的rootfs文件系统，Linux启动过程中将initramfs的内容解压到rootfs中。作为理解，可以认为initramfs挂载成Linux的临时根目录。}。内核初始化完毕后直接执行initramfs中的/sbin/init而不理会root=命令行参数。由initramfs中的/sbin/init程序负责加载磁盘驱动，查找根分区，挂载根分区到 /sysroot 并调用pivot\_root()\footnote{使用命令　man 2 pivot\_root查看介绍。}将找到的根目录替换为真正的根目录。然后执行真正的/sbin/init程序。

initramfs是一个目录的压缩包。由dracut程序生成。initramfs用于寻找真正的根目录时的临时根目录，通常只包含一些必要的驱动和一些用于查找根分区的脚本，但有时候可以成为一个完成的系统的根目录，这样的系统可以脱离磁盘运行。比如Fedora的安装DVD,使用一个巨大initramfs作为安装环境的根目录。该initramfs包含的是一个真正的OS，带有完整的工具。


\section{汉化man手册}\label{FAQ:zhman}

man使用程序使用nroff程序格式化man手册，然后使用less向屏幕分页输出。为什么呢？因为man手册的显示会随着终端的宽度不同而变化。所以man手册并不是使用原始的ascii文本文件编写的。
man手册源文件使用非常隐晦难懂的一些标记，这些标记借助nroff程序格式化后就获得了适合特定宽度的终端显示的文本。正因为如此，所以man手册能在各种大小的终端上显示而不破坏风格。

GNU使用的nroff是groff（GNU出品，带个g是常有的事情），但是groff并不支持UTF-8编码的文本。所以导致man手册中不能出现非英文字符。所以有人对groff进行了修改，添加了UTF-8支持，这就是groff-utf8。要想显示中文man手册，按照man-pages-zh\_CN这个中文手册之外，还得安装groff-utf8并修改/etc/man.conf，让man调用groff-utf8而不是原来的groff进行格式化操作。

\begin{code}
\#emerge man-pages-zh\_CN groff-utf8
\end{code}

编辑： /etc/man.conf 

找到这一行

{
\tt
NROFF           /usr/bin/groff -mandoc
}

替换为

{ \tt
NROFF           /usr/bin/groff-utf8 -Tutf8 -c -mandoc
}

然后一些汉化了的手册就不会再显示乱码了。
gentoo-zh overlay 已经集成了这一变化，groff-utf8 安装完成后会自动修改/etc/man.conf，如果修改没完成可以手工进行。

\section{符号链接}\label{faq:symlink}

从需求上看，有时候我们希望能使用不同的文件名访问同一个文件。符号链接就是干这个活的。可以理解成一个快捷方式。符合链接分软链接和硬链接两种。硬链接只能用在文件上而无法对目录使用硬链接。而且只有在同一个分区上才能使用硬链接。软链接没有这方面的限制。

硬链接就是在文件系统上出现了使用同一个磁盘区域的多个文件。所以会有上述限制也不奇怪。软链接实际上是个特殊的文件，文件的内容是链接的目标的文件名。打开软链接，实际上会打开链接到的文件。使用readlink(2)系统调用才可以获得链接目标的文件名。

符号链接使用 ln(1) 工具创建。软链接要加“-s”参数。

创建链接的语法是

\begin{code}
ln [-s] 链接到的目标 链接名称
\end{code}


“链接到的目标”可以使用相对路径。如果是相对路径，并不是相对执行命令的时的当前目录，而是相对于“链接名称”的路径。

\section{内存映射IO-MMIO}\label{FAQ:MMIO}

CPU对外设的操作有两种形式：专门的 IO 指令和 MMIO。 和内存拥有地址意义，外设也占有IO地址。IO地址和内存地址是独立编排的。对内存的操作使用的都是一般的内存操作指令，如 mov 。
而对IO地址的操作使用的则是专用指令 IN 和 OUT。IO指令是特权指令，只有CPU处于特权级别的时候（执行内核代码的时候）才能使用。

除了专用的IO地址和IO指令，CPU对外设还可以采取内存映射的IO。在“内存控制器”\footnote{注意和MMU的区别}的配合下，将一些PCI设备的IO地址映射到上内存地址。CPU即可使用一般的内存操作指令对IO进行操作。由于CPU对内存操作的指令非常丰富，这就方便了设备驱动的开发。当然，还有一个原因：x86 CPU只给IO地址分配了0--65535共64k个地址。而PC外设越来越多，IO地址显然不够用了，借用内存地址吧！

当然，也有部分外设是只能MMIO的，比如显存。显卡带有一部分直接被GPU访问的内存，这部分内存称为显存。CPU要访问显存，需要通过AGP/PCI-E总线，以操作AGP/PCI-E控制器的办法间接访问显存，显然这样会太麻烦。不如直接采取MMIO的办法，把显存映射到某个物理地址区域，以对待内存那样对待显存。对CPU来说，显存不过是访问速度慢了点的内存。但是这带来了一致性。

\section{内核视频模式设置-KMS}\label{FAQ:KMS}

\section{CISC和RISC}\label{FAQ:CISCandRISC}

在开发CPU的人群中，诞生了两个派系：CISC和RISC。（人类从来不缺乏派系斗争啊！汗！）
CISC是{}Complex Instruction Set Computing{}的缩写。RISC全称为{}Reduced Instruction Set Computing{}。

CISC{}将指令设计为单条指令可以执行多个低级操作。如一条add指令即可完成访问内存获取数字，并和寄存器中的数字相加这两个操作。
CISC的指令灵活性比较强，完成同样的操作可以使用多种指令的组合方式做到。而且每种组合执行效率都不一样。
RISC则反过来，每个指令只提供一种低级操作，CISC中的典型指令通常需要多条等价的RISC指令。但是由于RISC指令比较简单，每条指令所做的操作非常简单，所以多数指令可以做到单周期执行。
而CISC由于指令比较复杂，使得译码电路也大大的复杂化。执行每条指令所需要的时间也不相同。

简单的说，每条CISC指令其实可以分解为多个更简单的微操作，而每条RISC指令就是一条微操作。通常无法继续分解。另外，多数RISC指令都设计为等长格式，不论何种操作的指令，一条指令的
长度都是固定的。CISC通常使用变长编码。指令依据所要进行的操作不同而占据不同长度。

\section{RPC}\label{FAQ:RPC}

在计算机中，远程过程调用指的是程序调用一个子过程，子过程在另一个进程（可以是本机的另一个进程，也可以是网络上的另一台主机上的进程）里执行，但是编程上不需要关心网络传输细节。
只需要执行一次普通的函数调用即可。这样的过程就是远过程。远过程调用实现上有很多。桌面上常用的是DBus。Web编程有XML RPC 和 JSON RPC等可选。

\section{INI文件格式}\label{FAQ:INI}

INI 格式是 微软在 MS-DOS 开始沿用到 Windows ME 的系统配置文件的格式，通常使用扩展名.ini 因此得名。虽然后来被“注册表”所取代，但是并没有完全消亡。INI 文件因为其简洁的风格，被大量的软件用作其配置文件格式。

INI 文件的设置以 键值对为基础。如

\begin{verbatim}
key=value
\end{verbatim}

同时 ini 文件可以将这些键值对分节，节名用 中括号表示，如

\begin{verbatim}
[section1]
key1=value1
[section2]
key2=value2
\end{verbatim}

这样的文件就有2个节。key2 属于 节 section2 . 不同的节下允许同名的 key 。


另外 ini 文件也支持 以 “;” 号起始的行作为注释。如


\begin{verbatim}
[section1]
key1=value1
[section2]
; 注释哦
key2=value2
\end{verbatim}



\chapter{ 名称缩写	}
\chapter{ 图片索引	}


\chapter*{附录 \Alph{chapter}  Gentoo 的前身今世}
\addtocounter{chapter}{1}

\begin{notice}

声明: 本文的原始版本最初发表于IBM developerWorks，现在所有权归属Westtech Information Services。

本文档的中文翻译由\textbf{范华}完成。经范华同意，将本文以附录形式列出。

\end{notice}

\vskip 10ex

{\huge\bf Gentoo的前世今生，第一部 }
\addcontentsline{toc}{chapter}{Gentoo的前世今生，第一部}
\vskip 10ex


\section{Gentoo Linux发行版的诞生}

\subsection{我和linux}

现今对每一个linux爱好者（linux geek）来说，linux不再只是一个字面上的名称，她所呈现的一切对很多开发人员来说已经超过了他们所接触过的任何东西，linux比它们更强大、更令人着迷和称赞。当我在新墨西哥大学担任系统管理员时便与linux结下了不解之缘。那时因为我们的NT服务器运行得非常棒，我的手头上也有了很多空余的时间可以加以利用，就这样我的第一个linux操作系统（Debian）在一台Pentium 166的主机服务器上安下了家，接下来的不断、不断地学习和深入理解的过程使我对linux越来越着迷了……

一开始我学习了linux下的很多细节的东西：网络访问、执行备份、搞定samba等等。接着我建了一个qmail和apache的服务器并学习了python编程和shell编程。接着我还搭建了一个小型局域网接着就把linux娶回了家，在尝试过很多发行版后我最终选择了Stampede Linux这个版本\footnote{该版本从2001起就没有再更新了。}。你知道学习linux的过程是怎么样的吗？

\begin{itemize}
\item[]  第一: 努力搞清楚linux基本的东西
\item[] 第二: 当你已经有了相当好的掌握程度之后，学习定制你的linux，知识的累积会和你深入的程度成正比。
\end{itemize}

由于linux并没有隐藏任何东西，当linux对你来说变得越来越得心应手之后就可以开始探究技术和那些实现这些技术的工具了。


\subsection{Linux的潜能}

Linux提供了很多以前我所没有见到过的东西，如果一定要我用一个词来形容这些不可思议的话，我选择“潜能”这个单词：用来维护、改变、提高事物的能力，这种能力甚至能够冲破一些固有规则的约束。当我把kernel升级到一个更新的版本时，简简单单的就把我眼前的这个linux的性能提升了很多，更为令人兴奋的是这种改变几乎每时每刻都在进行着。而我也正是这种进步的一份子，伴随着linux的前进而不断进步着，对我而言这种感觉真的很棒。

如果你和我是同一类人，在你进入开源世界和linux世界之前大概看过位于Redmond和Cupertino的那些大公司们准备的下一代操作系统，它们确实如你所愿般的完美，然而那些东西却始终都只是一个虚幻的影子而已。然后就在我们慢慢等待的过程中linux来到了我们面前。虽然等来的这个精灵并不如我们预料的那么完美，但是她却提供给了我们这些喜欢动手hack的帅哥和美女一个亲手改变她的机会。就这样我们一边期待着一个更强大的操作系统，一边津津有味的hack我们的linux。日子一天一天过去，直到某天我们才突然发现原来期待着的那个强大的操作系统其实就在我们自己的手中，大家不约而同的笑了起来，也决定了继续在linux这条路上走下去，并继续hack下去。

\subsection{Linux的人文艺术}

我学到的另一件事就是Linux对人们的影响，这个话题可能听上去还真有点新鲜，是吧？Linux不仅仅只是一堆源代码的，它其实就是一个“社区”，从一开始的依赖这个社区解决我们提出的问题到付出我们的时间和经验去帮助他人，渐渐的我们也将会成为了这个社区的一部分。

IRC（Internet relay chat）既是一个交朋友的好地方也是一个很打发时间的场所。irc.openprojects.net上的\#stampede频道已经成为了我在网络上正式的安乐窝
\verb ^_^  。

那是我解答自己疑问的地方，也是第一次回答朋友问题的地方。\#stampede频道需要很多有安装经验的用户去帮助那些新手解决他们刚刚开始安装后碰到的各种各样的问题。由于那些新手在安装过程遇到的问题在irc中越来越普遍，原来很多有经验的Stampede Linux用户渐渐失去了他们一开始的热情。但是我依然还是很兴奋，因为很多菜鸟的问题我都知道解决的办法，要我忍着不回答那些问题我可做不到！当然我也并不是唯一的那个对解决新手问题乐此不彼的人，同样的家伙也有不少。我也承认自己也有那么点私心，想从那些更有经验的家伙们（不是指Stampede的开发人员）身上学到更多的东西。

\subsection{如何起步}

当有朋友问我如何才能加入一个开源项目时，我告诉他们的是首先是找一个能为他人做些什么的地方，就算那里只是解答一些很基础的问题。一份诚挚的渴望帮助他人的愿望是通往Linux社区的通行证，因为这份诚挚的愿望同样也扎根在每一个开源项目开发人员的心中（不仅仅只对Linux项目来说），它也应该扎根在那里。

沿着这条路走下去不可避免的你会遇到比你更有经验的同志，你将会从他们身上学到更多的知识，就像以前新手从你身上学习时一样。另一方面，当你积累起了更多的经验后在碰到某些问题时你就会用一个新方法去解决它而不是用以前惯用的一套思路。你遇到的一些开发人员有时会提出一些建议，有时又或者会需要一些帮助，他们更可能会邀请你加入他们的开发队伍；如果你的助人为乐成为焦点时，他们可能会笑着从你身边经过；如果你帮助了很多很多人之后，你在社区内肯定会备受瞩目。在Stampede和我身上这些故事都曾经发生过。

渐渐的我在Stampede的开发越来越深入，不久以后我就成为了一个正式的Stampede开发人员。在受到了Stampede的领导者Matt Wood的鼓励后，我开始对用于Stampede Linux软件包的原有的.slp机制进行升级。当时，.slp软件包格式包含一个.tar.bz2的软件包和后面的一个包含软件描述及软件包创作者等等在内的一个定长的页脚。这种实现的方式有两个主要问题：页脚部分实际上包含的内容根本达不到定长所约定的字节数；该格式没有预留任何扩充余地（也就是说如果未来没有办法加入一些可能需要的额外信息）。显然这些问题需要动一次大手术了，hoho。

和那些老资格的Stampede开发人员工作一段时间后，我拟了一个解决上面那些问题的草案。过了一阵子我便开始用Python先编写了一些原始的实现方案，新的格式(代号slpv6）有些类似与Amiga世界的IFF格式。下一代的.slp格式包含了了2的32次方个字段，字段种类为2的32次方种，每个字段最大数据段同样为2的32次方字节。新的格式不仅具有良好的扩充性而且比纯文本更加紧凑和简洁并易于解析。二进制代码和文本都能存储在这样的格式当中，该架构对其本身在未来的进一步发展带来了无限的可能性。我的想法是把这个新版的动态header加入到打包文件的结尾部分，从而这个新版本的.slp格式未来可以为Stempede用户服务相当一段时间并且同时又能和标准的UNIX档案文件保持不错的兼容性。

\subsection{丑陋的一面}

slpv6的开发进展很顺利，所有的资深开发者看到我取得的成果后都很高兴。不幸的是，两名刚加入的Stampede开发者想要自己掌控slpv6项目。由于不欣赏我选择的开发方向，他们花了很大劲诋毁和打击这个新的slpv6系统，虽然我也用了大量时间一边继续我的开发一边加入讨论再一边回应他们的攻击，但是这样做也没从根本上解决问题。最后一切都变的很明了，他们只是很擅长辩论，并且显而易见的是除非是走他们自己的路子，不然他们是不会罢休的。幸运的是我的项目依然得到了资深开发人员的认可和支持。可是这些讨论渐渐地使我背上了一些包袱，同时对Stampede的开发也产生了一些不好地影响。唉……

可惜我没办法使这些家伙消失，原来还可以在\#stampede频道里和那些高级的开发者互相交谈，但是现在不得不退了出来。每次只要我一进入那个频道，他们就开始变得很不友好，总是在破坏我想要进行的工作。这些家伙会使用各种各样的方法：比如一个开发者会议（其实只是想当着其他资深开发者的面侮辱我）。他们还尝试用投票的方法控制Stempede，当然那种投票只在他们可以得到更多支持的时候才会举行。但是自始至终我在这样的情况下都没有放弃过我的slpv6的开发工作。不用多说，资深开发者们都喜欢我的开发项目，并且他们也都支持我继续做下去（没有他们的支持，我不可能克服那么多困难并坚持下去）。

\subsection{对这些异类的了解}

我习惯于把这两个家伙和这种类型的开发者称为“异类”。虽然我的开发工作因此变得很很不愉快，但是我还是学会了怎么样去对付他们。就这点我乐于提供给各位一个对这些“异类”的全方面的描绘：他们的品质、采用的方法以及当你作为一个项目领导者怎么样才能对抗这些”异类“或是尽可能的用最小的代价去改变他们。

为了消除情绪上可能存在的危险，你需要具备一个先决条件：意志力。如果你不能用一种既礼貌又态度坚决的方式回应你的对手，事情就会变得很糟糕。“异类”的目的就是尽可能多的在你的项目中取得控制权，这么做会使他或她感觉更具有力量。首先，他们会对某个项目或是项目的开发人员进行片面的指责和抱怨，同时他们也会阻挠那些对这个项目富有建设性的提议。当然这些家伙在他们获得项目管理人员位置之前是不会对这个项目伸出任何的援手。目的就是使你确信只有依靠他们的那些“独道的、富有素养”的眼光才能最终解决问题，这样你就不得不给他们足够的权限去实现这些。

如果指责和抱怨没起什么作用，这些“异类”就会要求举行一个开发者会议。这将会给他们一个可以分裂你开发团队的机会。在觉得本方这方面已经得到了大多数人的支持后，他们就会举行一次投票决定（当然他们知道赢的会是他们的情况下）。如果并没有赢得投票或是投票被驳回，那么下周他们还是会提出举行一次会议以便再一次的分裂你的团队，然后再是那种无休止的循环。

如果会议的方法行不通，“异类”们将会变成革新运动者。他们会用一种更民主（也就是更容易操纵）的办法来取代先前压迫性的和非公平的决策方案。这些办法常常包括令人信服的让你去为你的开发团队中的大部分人做任何事。异类比较偏爱这个办法，因为你没有办法弃大多数投票表决的结果于不顾（活活活）。你许可这些事情发生的时候就已经把那把通往你的“Lexus”的“钥匙”交到了他们的手里，这将使你感觉越来越无力。

“异类”们用的另一种方法是激怒你的主要开发人员并使他们离开，然后在你的开发团队混乱的时候尝试重新组织该项目的管理团队。如果所有的努力都没有成功的话，他们会聚集尽可能多的叛离者并把他们安插在你的项目中，痛啊！

\subsection{对付这些异类}

区分这些家伙还是相当容易的。他们不会写一行代码（也不愿意写），相反他们会花大量的时间讨论那些更重要的问题（对了，就是那些管理方面的问题）。
假设你是一个项目管理者，对付他们非常容易。只需要告诉他们，在没有看到高质量的代码之前你是不会考虑他们所谓的建议的。
或者在他们提出“建设性“的批评之前强调对于某个项目有建设性得帮助也包括服从项目的管理人员。
如果他们开始编制优质的代码并且越来越有易于这个项目，那么就太好了。如果没有，就告诫他们离开。
在你忽略这帮家伙一段时间后，他们会选择离开或是一边采取行动一边写一些代码，世界就这样清净了\verb ^_^ 。

不幸的是Stampede的那些资深开发人员对”异类“并没有采取更多的管理措施。换句话说，他们许可了这两个家伙对我（和其他人）的无休止的纠缠。虽然这些资深开发者总是赞赏我的项目，但是对那两个家伙他们却并没有做的更多。然后终于有一天我决定制作一个自己的发行版，因为我觉得这样做比忍受那两个家伙更容易些。我退出了Stampede的开发团队并开始制定自己发行版的一些计划和草案。

一段时间之内，我对自己因为两个低等级开发者而离开一个项目还是感到有些不可思议。其实他们没有涉及到的实际情况却真正显示出这个项目存在很严重的管理方面的问题，如果高等级的开发人员不能或者不愿意确认Stampede的开发成果是可喜的和有益的话，我想我不会愿意继续留在那里。

\subsection{新的开始}

离开Stampede后我做的第一件事就是长长的舒了口气。喔……，整个世界都清净了。现在我有了足够的时间来思考我自己的Linux发行版的轮廓和将给Linux发行版的布局带来什么新的贡献。对Stampede感兴趣的一件事是它所具有的原生的性能（这得感谢它使用的带有实验性质的、并针对Pentium处理器优化过的pgcc编译器），所以我决定首先我考虑的就是性能。除了更少的CPU占用率以外，我还希望它更精简。很多发行版本（特别是那些流行的热缩塑料封装的家伙）默认启动了太多的daemons以至于打开一个xterm（X环境下的终端）后系统所剩余的可用RAM已经所剩无几了。我希望自己的发行版能更小也更强，为此我把目光放到了最大限度的榨取让这个操作系统运行的硬件平台的性能上。为此我下决心进行一个整体测试并处理掉所有细节中的性能方面的问题。

但是我真的很缺乏对应的资源，因为我是这个发行版的唯一的一个开发人员！我该怎样做才能只靠自己就鼓捣出不逊色于Redhat或是Caldera这样的产品呢？解决办法是采用自动控制技术。我必须写一些脚本以便所有的事情都可以自动搞定，这样我就可以事半功倍了。毕竟，电脑们这些方面做得更好，对吧？

很快我发现光是写一些自动化的脚本还远远不够，需要设计的是一整套能从源代码产生一个完整Linux系统的机制。我实验性的把它称做ebuild系统并且开始了工作。ebuild系统可以自动的建立所有一个发行版所需要的二进制文件，包括从解压源代码并打好相应的patch再到编译、封包的一系列过程的自动化解决方案。在一个基本、原始的ebuild可以工作后，我开始为一个Linux发行版必要的一些关键组成部分（像是gcc、glibc、binutils、util-linux等软件）撰写ebuild脚本。通过重新撰写初始化脚本（基于以前我为Stampede设计的初始化脚本）把原先的Stampede开发系统逐渐的演变成一个我自己的系统，接着我开始测试每一个自己建立好的新的软件包。

几个月之后我有了一个完整的，自主的Linux版本。我给她起了个名字『Enoch』然后坐着满足的笑了起来。但是什么改变了Enoch、Gentoo的发展又是怎么样的？续篇将会告诉大家Enoch是怎么演变成Gentoo的，而我在这条路上将要面对的更多新的挑战。

\vskip 10ex
{\huge\bf Gentoo的前世今生，第二部 }
\addcontentsline{toc}{chapter}{Gentoo的前世今生，第二部}
\vskip 10ex

\section{从Echo到Gentoo}

\subsection{Enoch踏出的第一步}

我在先前的文章中告诉了大家那段和Stampede开发团队在一起的、曾经最兴旺的时光和最后为什么离开的原因（就是想离那些有低级政治目的的、想控制项目的那帮家伙远点）。因为这些爱管闲事的好事者的干涉，我才会觉得装配一个自己的Linux发行版比在那种恶劣条件下改进Stampede要简单的多。幸运的是，我离开Stampede时是带着满满当当的经验离开的，这些经验与在Stampede的工作（应该是实质性的吧？）是分不开的，维护一些软件包也好、设计初始化脚本也好或是领导slpv6（下一代软件包管理系统）都使我相关方面的知识和经验得到了极大的丰富。

Enoch是我开始工作的这个版本的一个代号，得益于为它开发的高智能的包管理和升级系统，它将会是一个速度飞快的版本。我不得不承认这套智能化的系统在整个版本中占据了很大一部分位置，因为对于我这个光杆司令来说在那种重复性的劳动中消耗时间是没法接受的，所以才会要求开发中的系统必须自动为我完成那些琐事。另一方面完全由源代码来构建整个发行版（比那些“spin off”的版本、例如RedHat要好）也需要把工作划分好并尽可能多的挤出空闲时间来做这些工做。

使最基本的Enoch系统启动和运行之后，我回到了irc.openprojects.net并开设了自己的＃enoch频道。在那里我逐渐聚集起了10个开发人员组成的团队。在早期的那段时间里我们整天都聚集在IRC里，用空下来的时间制作我们的发行版。在我们无私的付出和大家的齐心协力的hack下，在不断的消除bug和新的bug的过程中，Enoch每天都在变化着，不管是专业化的程度还是各方面的功能都变得越来越出色。

\subsection{Enoch的第一块绊脚石}

不可避免的一天，Enoch碰到了它的第一块绊脚石。在加入了Xfree86、glib、gtk＋之后，我决定把xmms（一个基于X11/gtk＋的MP3/CD播放软件）弄进我的发行版，因为也该到了用音乐来调剂调剂的时候了！但是在安装好xmms之后启动它时......X死锁了！最初我觉得是自己使用的编译器的优化参数造成的（"-O6 -mpentiumpro"，在你看来有点诧异吧？）。第一个想到的解决办法就是用标准的编译器选项来编译，但是问题依然没有解决。然后只好到处寻找解决方法，接下来整整几个星期的开发时间我都用来追踪这个错误。一天，我收到了一个叫Omegardan的Enoch使用者的电子邮件，他也同样碰到了xmms的这个死锁问题。

交流了一段时间然后历经了n个小时的检测后我们发现死锁的原因在于POSIX的线程描述符（POSIX threads-related issue）。因为一些原因，pthread\_mutex\_trylock()函数没有返回它应该返回的值。作为一个Linux版本的创始者，这种类型的bug是我真的不愿意碰见的家伙。我指望开发人员能能够释出完美的源代码以便我可以把精力放到提高Linux易用性上，而不是把时间花在修复别人源代码的bug上。当然很快我就发现这种希望仅仅只是一个美好的想法罢了，相同的错误有时还是会出现。

在找到问题后，我们发现它不是xmms本身的问题，也不时gtk＋或glib的问题，也不是Xfree86 3.3.5没有thread-safe和死锁的问题，而是令人惊异的存在于Linux 的POSIX的线程执行本身，具体来说就是版本2.1.2的GNU C库（glibc）的部分代码中存在bug。我很震惊的是在Linux如此核心的部分居然存在这样严重的bug（而且我们为Enoch使用的glibc的版本是它的release版本，并不是什么prerelease版本或是CVS版本！）。

那么怎么样才能解决这个问题呢？我们不可能马上就能拿出一个修补方案，但是在浏览了一堆glibc开发人员的邮件列表后，我偶然发现了还有一个人也碰到了相同的问题，然后在glibc开发人员在回复他的邮件里我们找到了那个附带的补丁，它为我们解决了那个线程问题。但我令我好奇的是为什么同样使用glibc 2.1.2的RedHat 6没有受这个bug的影响（当时RedHat 6的发布时间先于那个补丁的出现）。为了找到答案，我下载了RedHat里glibc的SRPM包（source RPM）想看一下他们使用的补丁是怎么样的。

RedHat有他们自己的glibc补丁来解决pthread\_mutex\_trylock()函数的问题。显而易见的是他们也碰到了同样的问题，然后自己进行了修补。但是由于RedHat没有把这个补丁回馈到glibc的开发社区，其他人们就没有办法分享这个补丁。但是也可能是RedHat把这个修补方案回馈到了glibc的开发社区，然儿glibc的开发人员并没有接受这个修补方案。 或者这个bug只会在特定版本的binutils和特定版本的编译器连用时才会触发，然而RedHat使用的binutils和编译器的版本并不是这两个特定的版本（虽然RedHat还是给出了这个补丁）。我猜测我们永远也不会知道究竟事情的真相是什么样的，但是我学会的一件事情是：RedHat的SRPM包里有很多定制的补丁和增强代码，而这些代码和补丁看来从来没有回馈到原始的开发人员那里。我将会为此来上一段激昂的演说。

\subsection{激情的演说}

当你将一大堆各种各样的源代码汇聚成一个Linux发行版时，把所有你做好的bug fix和补丁反馈给原始的某个软件包的开发人员是一件相当重要的事情，就如我了解到的那样，这是发行版的开发人员为Linux做贡献的很多途径中的一个。我们也恰好就是这样的一群人，为的就是把很多不同的程序和软件集合在一起，让它们工作起来就像是一个整体。将来我们也会把我们们对一些软件所做的修改和补丁反馈回原始软件的开发人员以便其他的用户和后来的发行版能从中受益。如果你只是把补丁留在你自己那里，这样做不会对任何人有什么帮助，很多人们将会为一些相同的问题浪费掉大量的时间。这种不顾别人的方式违背了整个开源世界的精神和宗旨，同时对Linux的发展也只是有害无益。或许我应该说这样的做法对我们来说就是一个大大的“BUG”。

不幸的是一些发行版（啊咳）（RedHat）并不如其他一些版本（Debian）那样对整个开源社区分享他们的成果。

\subsection{编译器的艺术}

在我们尝试解决glibc 线程问题的时候，我给Ulrich Drepper发了封email（他是Cygnus的一员并且在glibc的开发中举足轻重）。我在e-mail中提到了我们碰到的POSIX线程问题和我们在Enoch中使用pgcc来获得优化的性能。在他的回信中他这样提到（我解释一下）：“我们自己的包含在CodeFusion中的编译器制作的可执行代码比其他的一些编译器、比如pgcc编译出来的代码执行速度更快速。”显然我对测试测试Cygnus那帮家伙开发的神秘的“turbo”编译器非常有兴趣。

因此我申请拿到了一个Cygnus Codefusion 1.0的demo拷贝以便我可以对它的性能做一个测试。Omegadan和我对测试的结果很吃惊，它同Ulrich提到的那样出色。x86的后端提高了90％的有关cpu-intersive的可执行文件的执行效率（比如bzip2）。几乎每一个程序都能从中获得至少10％的真实世界的性能提升，而我们所作的仅仅是换了一个编译器。Enoch的速度也因此获得了30％－40％的提升。同时性能也提高了不少，提升的幅度超过了我们以前把编译器从gcc切换到pgcc时提高的幅度。显然，在对这个编译器的测试后，我们很希望把这个编译器包含在Enoch中，有点幸运的是CodeFusion CD中的包含的源代码遵循的是GPL，这样在Enoch中使用这个编译器已经可以算是已经得到了完全的认可了..........，至少我们是这么想的。

\subsection{异常事件的发生}

为了能在Enoch中使用这个编译器，我给Cygnus的市场部主管发了一封电子邮件，但是期望中的“哦，拿去用好了，感谢使用我们的编译器！”这样的回复并没有收到，取而代之的是一句“虽然在技术上我们许可使用Cygnus的编译器，但是我们强烈建议不要在在Enoch中使用该编译器或是包含它的源代码。接着在我的回复中我问了他们这样一个问题：“既然不愿意让别人使用它的源代码，为什么还在以GPL的许可条例来发布它的源代码？”作为一个猜测，我觉得他们事实上是不想以GPL的方式来发布他们的源代码的，但是由于这个编译器是源自egcs（以GPL方式发布的），他们除了以GPL方式发布之外别无选择。

这是当某一个公司想使用开源的代码来生产私有产品这样的情况时，GPL如何阻止这样的事情发生的一个很好的例子。我比较有根据的一个猜测是Cygnus担心我们使用这个编译器后将会打击到他们整个产品框架的销售，更加奇怪的是不管是他们的行销方案还是InfoWorld的预览中都没有提及包含在CodeFusion中的那个新的编译器，因为CodeFusion销售的是一套“development IDE”而不是一个编译器。

为了缓解一下他们那种偏执的态度，我提出了个建议，就是在我们的Enoch主页上放置上CodeFusion的签注文件并加上一个链接来刺激CodeFusion的销售。从我个人的观点来说，我不认为一个“turbo”的Enoch会影响到CodeFusion（虽然它是一个IDE产品）的销售情况。但是我还在想方设法的令到他们愉快，比如告诉他们这个IDE的组件是一个商业化的产品，我们也并没希望或者有什么意图用Enoch来发行它。

我把这个（大方的）请求用电子邮件的方式发给了Cygnus，但是收到的确实另一个奇怪的回复。他们想得到所有我们关于“市场元素”方面的具有权威的权利（显然，这也包括了我们网站上的内容），真是太令人震惊了。Cyguns的营销团队似乎对Linux社区和GPL的运作一无所知，事到如今我终于决定终止与Cygnus彼此间的联系，因为再这样下去事情会变得怎么样谁都不知道。与此同时，我们为Enoch准备了两个版本，一个是内部的“turbo”版，一个是公开的“non－turbo”版，其实就是把决定留在将来再去做。

但是几个月之后，他们就把CodeFusion x86的backend换成了gcc 2.95.2，现在不只是那些知道包含在CodeFusion CD中的“隐秘的GPL编译器”的这群人可以获益，几乎每一个人都可以从这个新的优秀的backend中获益了。然后我们还是决定继续前行，尽量使用gcc来替代CodeFusion的编译器。在gcc 2.95.2已经越来越成熟的情况下，我们已经可以放开Cygnus了（同时，RedHat却为购买这个CodeFusion而花费了比较冤的一笔钱了。）（注：新的x86版本gcc 2.95.2的backend为新的Linux发行版提供了一开始我们提到的很重要的速度提升，它也为FreeBSD 4.0相对3.3.6版本速度上提升做出了很大的贡献。你注意到这两个提升的不同点吗？）

\subsection{肥皂盒}

感谢这件事情和其他的一些经验，我从中对那些以开源为主要获利手段的企业有了很深的理解。虽然对个人来说，乐于生产私有闭源软件这件事并没有任何错误的地方，但是一个开源企业搅乱或是拒绝与其他的开源世界合作是没有任何意义的；同样，不支持GPL或是其他的等等也没有什么意义。这是一个实践性质的并具有现实意义的观点。

思想和代码上自由的交换才是开源企业得以获利的根本，这点他们应该充分的认识到。反过来，对立与GPL标准只会破坏这个他们依赖于发展与繁荣的环境。换句话说，开源的环境是你事业的土壤，保护这片土壤的纯净还是很有意义的。

我也懂得在短时期内保留一些代码上秘密来获得财富的累积是一个颇具诱惑性的东西，先进的代码和特别的技术提供给了人们一个在竞争中获得优势的绝好机会，由此可以获得增长的销售业绩和利益。但是当你的目的是成为一个唯一的产品提供者，而这个产品商业的成分大于开源的成分时，开源世界是不会许可这样排外性质地使用开源或是相关东西的，这就是开源的意义。

\subsection{回到Enoch}

现在，我从自己的肥皂盒中出来并继续我的故事。 

由于Enoch已经变得越来越出色，更名的计划也就这样列入了我们的议事日程当中，接着“Gentoo Linux”诞生了。然后就是朝Gentoo Linux 的1.0版本努力前进中。大约也是这个时候，我决定帮我那台Celeron 300M（超频到450M并且十分稳定）的老电脑升级一下，新平台是一块崭新的Abit BP6主板（从市场上找到的双Celeron接口的）。在卖掉了老主板后我把我两个Celeron 366的系统集中起来，然后把Celeron 366超到了500Mhz，然后开始工作了。但是我注意到我的新机器不是非常稳定。 

显然我第一个反应就是把频率改回没超之前的366Mhz，但是随之而来却遇到了一个更奇怪的问题：不管CPU全速运转多少时间，系统都不会死锁；但是一旦空闲下来过一夜的话，系统有很大的可能就会完全死锁掉。是的，这是一个idle bug－－－－噢！在作了一些调查之后，我发现在这块主板上也有其他用户碰到了这个相同的问题。原因是BP6主板上的一个芯片（可能是PCI控制器）与标准规格有点不同或是比较古怪，这个东西就是造成Linux在空闲时候死锁的主要原因。

我渐渐的心烦意乱起来，因为我没法再去采购另外的PC部件了，Gentoo的开发也只好被迫终止下来。我也开始对Linux越来越有些悲观的情绪了并决定转向FreeBSD。是的，的确是FreeBSD！这部分就此为止了，我们第三部再见了：）

\vskip 10ex
{\huge\bf Gentoo的前世今生，第三部 }
\addcontentsline{toc}{chapter}{Gentoo的前世今生，第三部}
\vskip 10ex

\section{离开Linux后再度回归}

在前一篇文章的结尾部分，我说到因为新升级的双Celeron主板（Abit BP6）存在一个古怪的空闲时死锁的问题导致Gentoo开发停止。虽然解决问题的办法就是更换主板，但是我已经没有重新更换主板的资金了，这件事也打击了我对Linux的信心并使我决定中断Gentoo的开发并转向了FreeBSD。我需要的是一个可以正常运转的系统，而Linux在这个时候的表现并不尽如人意（一天到晚的死锁），那个当口，我觉得是好好接触接触FreeBSD的时候了，便在机器上安装了FreeBSD后开始了又一次的捣腾，在接下去的几个月中，我也几乎没有再碰过Linux一个指头。

\subsection{FreeBSD之印象}

首先，我真的很喜欢FreeBSD。我感觉这个操作系统是一个组合的很完美的系统，它的几乎每一个部分都同样精巧，而这种精巧的在Linux世界中几乎不存在。我的满意实质上是来源于那些FreeBSD中非常充足的man page，这可不像Linux里那些只有GNU info文档的很多软件那样让人根本没法用。

最最重要的是我对FreeBSD中维护与升级系统的ports系统印象非常深刻。与Linux维护与升级的方法不同，ports使用的不是二进制的软件包而是直接去原始的软件站点下载所需要的源代码并编译。不管你是安装Samba或是升级核心系统都是在你的机器上用源代码编译而成。这样的实现方法和我在Gentoo Linux中建立的那套机制有着异曲同工之处。从这点和其他许多方面来说，FreeBSD的这种设计符合我作为一个开发人员和一个系统管理员所期望的那种感觉。就这样，FreeBSD为我营造了整整几个月舒适的工作环境，同样我也很乐意于花些时间在这个出色的操纵系统中探求与获取知识。

\subsection{FreeBSD的优点}

很多Linux和FreeBSD之间的不同点都是源自与它们本身开发架构的不同。Linux的开发架构非常松散，我们只是依靠不同的发行版把分散在Internet上呈离散状态的很多部分组合成一个完整的Linux，而FreeBSD和其他BSD系统（OpenBSD和NetBSD）都有一个唯一的核心小组来确保源代码的单一性和协调性，这样至少每一种BSD自身都拥有一套统一的源代码设置。这是一件挺棒的事情，也是FreeBSD感觉上和Linux那种“patch集合”有所不同的主要原因。

接下来，我们在纯技术方面再作个比较。很多FreeBSD的粉丝都声称FreeBSD比Linux更合适用作服务器上跑的操作系统，他们会告许你在高负载情况下FreeBSD表现得更好，而且它的TCP/IP栈相对出色一些（如果你用Linux 2.2或更早版本的内核和FreeBSD作比较，我同意这个说法）。FreeBSD确实是一个很好的服务器操作系统，这点勿庸置疑，但是这只是FreeBSD相对Linux 2.2或更早的内核版本时的情况。我作为一个新版本内核的粉丝，早就在我的电脑上用上了2.4测试版的内核，它确是也很棒，从出色的TCP/IP栈到整个重新设计的“netfilter”系统都是。我觉得在不久的将来，新的性能标准将会由Linux来定义，而“free UNIX”将会在商业领域面对Linux强有力的挑战。

\subsection{FreeBSD的不足}

与服务器领域的应用不同，在桌面应用上，Linux占有绝对份额上的优势（仅相对BSD来说，Linux不管是对Win还是对MAC都完全处于下风）。所有最新的桌面应用软件一定是先在Linux上出现、在3D加速和声卡的支持方面，Linux也比BSD走在了前面。随着2.4版本内核的临近，Linux在这块地盘上还是会继续保持它的优势地位。

我对FreeBSD采用的UFS文件系统并不喜欢，虽然UFS相对Linux的ext2文件系统来说更健壮，但是付出的代价是那个另人昏昏欲睡的龟速。现在也有一个UFS文件系统的扩展叫“soft update”，它是把小块的IO操作聚合成大的文件块后再写入物理硬盘以提高文件系统的速度，就算“soft update”这套机制大幅提高了UFS文件系统的性能，我也没法就说在所有方面的比较中UFS都比ext2优秀。当然，UFS和“soft update”更加可靠，FreeBSD也可能会在文件系统的战争中击败Linux，但是请不要忘记，输给FreeBSD的仅仅只是现在的2.2版本或者更旧版本的Linux，这不代表将来也会。

现在，我们把话题转变一下，我们比较的双方是现今的Linux 2.2版本、2.4版本和FreeBSD。Reiserfs（一个新的日志型文件系统）已经给我们带来了一阵惊喜，而Linux还有蓄势待发的ext3、IBM的JFS和XFS文件系统，这些文件系统都在提供高可靠性的同时提供了优秀的性能。Reiserfs给了Linux在文件系统上超越FreeBSD的一个契机，这也是我认为Linux 2.4版本会上演大逆转的原因，FreeBSD的传统强项在未来2.4内核面前可能会荡然无存。

\subsection{回到Gentoo的开发}

几个月之后决定重新回到Linux世界的我在一台新的机器上又装了Gentoo。首先，回到Gentoo的开发中来是一个计算后的决定－－我已经花费了很多时间使自己成为一个Linux的万事通，而现在怀抱着BSD就等于是把以前学到的知识都浪费掉了，这样做我觉得不是很值得。而且在更新Gentoo Linux后那么一段很短的时间内，我为“为什么再次回到Linux怀抱”找到了几个新的理由，也就是前面提到过的kernel以及文件系统的改进等等。FreeBSD是一个宁静的家园，但是这样的宁静太安静了点，这样的宁静也包含着困惑。相反Linux世界充满着活力，发展也是日新月异。如果你所寻找的是兴奋和创新的地方，那么毫无疑问Linux就是你所向往的世外桃源。

Linux从2.0进步到2.2给我的感觉就是满失望的，但是2.4时代是绝对值得去守候着的，为此Gentoo Linux重新回到了我们面前，那种兴奋的感觉也重新回到了我的心中。

Gentoo Linux重生的另一个关键因素是我们开发团队的领导者－－Achim Gottinger。我想花一点篇幅对他所给予的帮助（使我我重新开始了Gentoo Linux的开发）致以诚挚的感谢。我在回到Linux世界之前就开始与Achim Gottinger有了电子邮件上的往来，在几乎每一封他的电子邮件中，我都可以看到一些新的.ebuild或者是些迫切需要修复的bug。在我回到Linux世界并重新开始了Gentoo的开发之后，Achim继续贡献着他的时间和精力使这个发行版步入正轨。直到最近，Achim和我都是Gentoo Linux仅有的两个开发者，这也是出于选择的结果。因为我们都使用几乎相同的发行版，也因为Achim的技术，我们可以轻松的完成非常巨大的工作量以至于我觉得加入第三名开发者并不会对我们的进展有什么帮助。现在Achim是Gentoo Linux开发组的负责人，几乎每天Gentoo的都会有基础部分中主要的提高。我们已经走到了这里，也已经准备好了CVS树为后来者提供一个协同开发平台，小心翼翼的逐步扩大Gentoo开发队伍的工作也开始付诸实施。

\subsection{新的版本}

我没有觉得花在BSD上的时间是在浪费。实际上，它给了我一个很好的机会来反省一下整个Linux社区存在的问题和Gentoo Linux应该做点什么来改进这些短处。. 

在新版本的Gentoo Linux中，我下决定不再使用pgcc或者什么非常优化的参数来编译所有的软件包，因为稳定性还是要放在第一位的，我们默认将会使用合理的优化选项（"-O2 -mpentium"），但也同时向用户提供了可以简单自定义的优化选项来满足了一些同胞希望得到最“bleed edge”的系统（通过我们的自动化系统完成）这么个愿望。

FreeBSD给了我一个关于“自动化定制系统如何工作？”这个问句一个很好的提示。我决定在我们的自动化定制系统（现在叫做Portage）中加入一些FreeBSD的特性来制作一个新一代的ports系统。 

Portage 可以说是Gentoo Linux的心脏，它所具备的东西远远超过一个简单的包管理机制或是一个系统管理机制。Portage通过它包含的对制作工具的设置和制作脚本可以使你从源代码构建一个完整的发行版系统。但对我来说更重要的是，Portage给用户提供了一个可以完全接触Gentoo Linux构建智慧的途径。对我们开发者来说，这意味着当Gentoo Linux不断发展的同时我们也记录下了一个发行版制作的过程。Portage的易用性和可读性也为越来越多的人提供了一个窥探Linux内部的窗口，它也为后来者贡献他们的代码和脚本打开了方便之门。

Portage是我们为他人展示Linux技术和原理的一条途径，通货学习自动化制作脚本，你可以看到大量各不相同的包是怎么互相适应并结合成一个整体的。如果你需要，你也可以从我们的站点上攫取整个CVS树然后自己hack并制作个人的Linux发行版。我们坚信这是一件好事情－－我们希望把知识交给渴望这些知识的人们以便他们可以把Linux带入一个新的领域。

\subsection{商业上的关注}

起初，有许多拥有不同背景的人们加入了Gentoo的开发中来。因为这个，我们的开发人员对于如何最终在Gentoo上获得经济利益也有许多各不相同的打算，对此我并没有太多的诧异。基本上有这么两种类型的开发人员：一类群体反对用Gentoo来追名逐利，另一类群体则对使Gentoo Linux成为一个成功的商业产品非常感兴趣。这是一个预料中会存在分歧的地方，第一类群体认为商业化的运作包含着腐化等不良的影响，而第二类群体则认为没有这么多的负面因素。

在以前还是Enoch的那段时光中，我对商业成份究竟有利还是有弊这点也很难做个了断。我验证过的是像Debian这样的Linux发行版真正忠于“自由”这样的事实，我喜欢这样。对比其他商业化的发行版，他们给用户带来的易用性包括了在各自的网站上提供一份完整的安装说明，这也是一个我想去借鉴的好东西。

同样，我也真心希望Gentoo Linux能够成为一个成功的商业版本，为了这个目的，我努力想在商业和开源之间找到一个平衡点，可是直到最近我还是没有能够找到这么一个黄金分割点。

\subsection{该做些什么}

我们该怎么做才能在商业化和非商业化中取得平衡呢？关键的一点是一定不能忘记我们的基楚和根本－－－Gentoo Linux 作为一个开源软件的根本和基础。所有我们作出的努力都必须遵循这个基础，这不仅仅是肯定开源软件或只是使用开源软件，还是对开源软件和开源发行版开发的鼓励和支持，也不会发对用这样的一个对待开源姿态来获取商业回报。更重要的是，我们绝不会采用商业化的模型，因为这样做对于其他发行版使用我们的源代码有阻碍作用。我们的开发团队对所有人来说都会是开放的和可接近的，而Gentoo Linux这个自由发行版不仅仅可以被大家接受还会因为很多人的鼓励而继续走下去。我们必会成为开源运动的倡导者，一个把这个理念贯彻到行动中而不是停留在文字层面上的倡导者。

如果某公司需要为一个商业化的基于Linux技术的需求使用Gentoo Linux，他们可以从我们的CVS树中攫取这些代码并马上开始使用它们，因为所有我们的分散的工作都是基于GPL。我们在确信所有基于Gentoo Linux的衍生产物都遵循GNU Public License的前提下是不会在任何地方限制别人使用我们的代码的。

我们希望有尽可能多的人们从我们的工作中受益，但是我们也希望尽可能多的能从你对Gentoo Linux的提高中获益。如果你公司的产品有很大一部份是基于Gentoo Linux的话，希望你可以把所有可分类的修改和提高发送给我们以便加入到CVS树中使更多的人获益。继续保管和改进你提交的修改后，你也能从我们所做的修改中受益。我们也鼓励商业实体和非商业实体之间的合作，举个例子来说：不管是在他的ISP中使用Gentoo Linux的系统管理员还是用Gentoo Linux构建商业服务器的公司都能从彼此对Gentoo Linux的改进中获益。是时候来促进在人们之间的自由代码交换了，这也只有开源软件可以做到。

\subsection{将来要走的路}

现在离Gentoo Linux 1.0 的发布已经很近了
（在你在developerWorks上读这篇文章的时候它可能已经发布了，想想现在的2006.0是不是大家有种沧海桑田的感觉
\verb ^_^ ??
）\footnote{作者在IBM DeveloperWorks上最初发表的时候，距离Gentoo Linux 1.0的发布还有时日}。但是Gentoo Linux将来的方向会是怎么样的呢？

当我们逐步迈向2.0版本时，我希望继续提升Portage作为Gentoo Linux核心的性能，因为任何关于Gentoo Linux主要的进步都会从Portage的进步开始。主要代码从bash转换到python的过程我也会继续下去，因为这么做会使我们加入新的设计（比如为我们的全自动构造系统设计的面向对象的新东东）。

除了Portage的修改，我还希望小心谨慎的寻找技术出色并且和我们使用相同版本的开发者加入我们的开发团队。在扩大了开发团队之后，我们可以为Gentoo Linux的加入更多的自动化定制脚本。比这更重要的是，适当扩大的开发团队可以使Gentoo Linux站在Linux技术的尖锋之上，这才是乐趣所在嘛：）

我们也希望商业化的Linux技术公司可以把Gentoo Linux作为他们产品的基础。现在我们已经有了这样一个关系，将来也会更多的，而这样的协作承诺充满着乐趣并对于Gentoo Linux的用户非常有益。

最后我要说的是，我们主要的目标是为Linux社区提供有意义的贡献。虽然可选择的发行版很多，但是Gentoo Linux还是拥有许多其他版本所没有的东西。我们对未来Gentoo Linux发展充满着信心，我们希望你也有同样的感觉。

\vfill

\begin{center}
{\Huge 全文完}
\end{center}

\vfill

\end{document}
